% Options for packages loaded elsewhere
\PassOptionsToPackage{unicode}{hyperref}
\PassOptionsToPackage{hyphens}{url}
\documentclass[
  british,
  a4paper,
]{article}
\usepackage{xcolor}
\usepackage[margin=1in]{geometry}
\usepackage{amsmath,amssymb}
\setcounter{secnumdepth}{-\maxdimen} % remove section numbering
\usepackage{iftex}
\ifPDFTeX
  \usepackage[T1]{fontenc}
  \usepackage[utf8]{inputenc}
  \usepackage{textcomp} % provide euro and other symbols
\else % if luatex or xetex
  \usepackage{unicode-math} % this also loads fontspec
  \defaultfontfeatures{Scale=MatchLowercase}
  \defaultfontfeatures[\rmfamily]{Ligatures=TeX,Scale=1}
\fi
\usepackage{lmodern}
\ifPDFTeX\else
  % xetex/luatex font selection
  \setmainfont[]{Noto Serif}
  \setmonofont[]{Noto Sans Mono}
\fi
% Use upquote if available, for straight quotes in verbatim environments
\IfFileExists{upquote.sty}{\usepackage{upquote}}{}
\IfFileExists{microtype.sty}{% use microtype if available
  \usepackage[]{microtype}
  \UseMicrotypeSet[protrusion]{basicmath} % disable protrusion for tt fonts
}{}
\makeatletter
\@ifundefined{KOMAClassName}{% if non-KOMA class
  \IfFileExists{parskip.sty}{%
    \usepackage{parskip}
  }{% else
    \setlength{\parindent}{0pt}
    \setlength{\parskip}{6pt plus 2pt minus 1pt}}
}{% if KOMA class
  \KOMAoptions{parskip=half}}
\makeatother
\usepackage{longtable,booktabs,array}
\newcounter{none} % for unnumbered tables
\usepackage{calc} % for calculating minipage widths
% Correct order of tables after \paragraph or \subparagraph
\usepackage{etoolbox}
\makeatletter
\patchcmd\longtable{\par}{\if@noskipsec\mbox{}\fi\par}{}{}
\makeatother
% Allow footnotes in longtable head/foot
\IfFileExists{footnotehyper.sty}{\usepackage{footnotehyper}}{\usepackage{footnote}}
\makesavenoteenv{longtable}
% definitions for citeproc citations
\NewDocumentCommand\citeproctext{}{}
\NewDocumentCommand\citeproc{mm}{%
  \begingroup\def\citeproctext{#2}\cite{#1}\endgroup}
\makeatletter
 % allow citations to break across lines
 \let\@cite@ofmt\@firstofone
 % avoid brackets around text for \cite:
 \def\@biblabel#1{}
 \def\@cite#1#2{{#1\if@tempswa , #2\fi}}
\makeatother
\newlength{\cslhangindent}
\setlength{\cslhangindent}{1.5em}
\newlength{\csllabelwidth}
\setlength{\csllabelwidth}{3em}
\newenvironment{CSLReferences}[2] % #1 hanging-indent, #2 entry-spacing
 {\begin{list}{}{%
  \setlength{\itemindent}{0pt}
  \setlength{\leftmargin}{0pt}
  \setlength{\parsep}{0pt}
  % turn on hanging indent if param 1 is 1
  \ifodd #1
   \setlength{\leftmargin}{\cslhangindent}
   \setlength{\itemindent}{-1\cslhangindent}
  \fi
  % set entry spacing
  \setlength{\itemsep}{#2\baselineskip}}}
 {\end{list}}
\usepackage{calc}
\newcommand{\CSLBlock}[1]{\hfill\break\parbox[t]{\linewidth}{\strut\ignorespaces#1\strut}}
\newcommand{\CSLLeftMargin}[1]{\parbox[t]{\csllabelwidth}{\strut#1\strut}}
\newcommand{\CSLRightInline}[1]{\parbox[t]{\linewidth - \csllabelwidth}{\strut#1\strut}}
\newcommand{\CSLIndent}[1]{\hspace{\cslhangindent}#1}
\ifLuaTeX
\usepackage[bidi=basic,shorthands=off]{babel}
\else
\usepackage[bidi=default,shorthands=off]{babel}
\fi
\ifPDFTeX
\else
\babelfont{rm}[]{Noto Serif}
\fi
\ifLuaTeX
  \usepackage{selnolig} % disable illegal ligatures
\fi
\setlength{\emergencystretch}{3em} % prevent overfull lines
\providecommand{\tightlist}{%
  \setlength{\itemsep}{0pt}\setlength{\parskip}{0pt}}
\usepackage{array}
\usepackage{tabularx}
\usepackage{bookmark}
\IfFileExists{xurl.sty}{\usepackage{xurl}}{} % add URL line breaks if available
\urlstyle{same}
\hypersetup{
  pdflang={en-GB},
  hidelinks,
  pdfcreator={LaTeX via pandoc}}

\author{}
\date{}

\begin{document}

\section{Regular-entry compression pilot
01}\label{regular-entry-compression-pilot-01}

\emph{Assembly-only pilot comparing the current full regular-entry
rendering with shorter book-like variants. Source model entries are
unchanged.}

\subsection{Sample}\label{sample}

This pilot samples 12 regular entries chosen to cover very simple nouns,
ordinary verbs, short note sections, dialect/source notes, and manual
comparison cases.

See \texttt{regular\_compression\_pilot\_manifest.tsv} for the selected
rows and sample reasons.

\subsection{Variant A: current style
baseline}\label{variant-a-current-style-baseline}

\subsubsection{bake --- OE bacan}\label{bake-oe-bacan}

Derivation: \emph{*bákaną} \textgreater{} \emph{bacan} (regular).

\paragraph{Derivation trace}\label{derivation-trace}

Proto input: \emph{*bákaną}

\begingroup
\setlength{\fboxsep}{6pt}

\noindent\fbox{%
\begin{minipage}{0.97\linewidth}
\small
\begin{tabularx}{\linewidth}{@{}>{\raggedright\arraybackslash}p{0.300\linewidth}>{\centering\arraybackslash}X>{\raggedright\arraybackslash}p{0.540\linewidth}@{}}
\begin{minipage}[t]{\linewidth}
\centering\textbf{Earlier Germanic changes}\par
\vspace{0.35em}
\raggedright
\centering\textbf{West Germanic}\par
\raggedright
\vspace{0.2em}
\raggedright [no change]\par
\vspace{0.6em}
\centering\textbf{Northwest Germanic}\par
\raggedright
\vspace{0.2em}
\raggedright [no change]\par
\end{minipage}
&
&
\begin{minipage}[t]{\linewidth}
\centering\textbf{Old English changes}\par
\vspace{0.35em}
\raggedright
\begin{tabular}{@{}>{\raggedright\arraybackslash}p{0.74\linewidth}@{\hspace{0.55em}}>{\raggedright\arraybackslash}p{0.16\linewidth}@{\hspace{0.25em}}}
Anglo Frisian Brightening & \emph{*bækaną} \\
OE A Restoration & \emph{*bakaną} \\
OE Heavy Syllable Nasal Apocope & \emph{*bakan} \\
OE Secondary Nasalization & \emph{*bakąn} \\
OE Weak Tail Reduction & \emph{*bakan} \\
\end{tabular}
\end{minipage}
\\
\end{tabularx}
\end{minipage}%
} \endgroup

Outcome: \emph{bacan}

\paragraph{Reconstruction and comparative
evidence}\label{reconstruction-and-comparative-evidence}

Orel reconstructs the verb as \emph{*bakanan} and cites Old English
\emph{bacan} beside Old High German \emph{backan, bahhan}
(\citeproc{ref-Orel2003}{Orel 2003}). Campbell gives \emph{bacan} as one
of the standard examples of Old English A-restoration before a single
consonant, and Ringe and Taylor state the same development from
\emph{*bakan} to Old English \emph{bacan}
(\citeproc{ref-Campbell1959}{Campbell 1959, 61};
\citeproc{ref-RingeTaylor2014}{Ringe and Taylor 2014}).

\paragraph{Old English evidence}\label{old-english-evidence}

Bosworth-Toller and Clark Hall both record \emph{bacan} as the ordinary
Old English verb `to bake' (\citeproc{ref-BosworthToller1898}{Bosworth
and Toller 1898, 72}; \citeproc{ref-ClarkHall1960}{Clark Hall 1960}).
The target in this entry is therefore the attested infinitive headword
itself, not a selected oblique or finite paradigm cell.

\paragraph{Development to Old English}\label{development-to-old-english}

From \emph{*bákaną}, Anglo-Frisian brightening first gives
\emph{*bækaną}. A-restoration then returns the stem vowel to \emph{a}
before single \emph{k} plus the back-vocalic infinitive suffix, and
later apocope and weak-tail reduction yield \emph{bacan}
(\citeproc{ref-Campbell1959}{Campbell 1959, 61};
\citeproc{ref-RingeTaylor2014}{Ringe and Taylor 2014}). The development
is therefore straightforward: \emph{*bákaną} \textgreater{} bacan.

\subsubsection{beech --- OE bōc}\label{beech-oe-bux14dc}

Derivation: \emph{*bōkō} \textgreater{} \emph{bōc} (regular).

\paragraph{Derivation trace}\label{derivation-trace-1}

Proto input: \emph{*bōkō}

\begingroup
\setlength{\fboxsep}{6pt}

\noindent\fbox{%
\begin{minipage}{0.97\linewidth}
\small
\begin{tabularx}{\linewidth}{@{}>{\raggedright\arraybackslash}p{0.460\linewidth}>{\centering\arraybackslash}X>{\raggedright\arraybackslash}p{0.460\linewidth}@{}}
\begin{minipage}[t]{\linewidth}
\centering\textbf{Earlier Germanic changes}\par
\vspace{0.35em}
\raggedright
\centering\textbf{West Germanic}\par
\raggedright
\vspace{0.2em}
\raggedright [no change]\par
\vspace{0.6em}
\centering\textbf{Northwest Germanic}\par
\raggedright
\vspace{0.2em}
\begin{tabular}{@{}>{\raggedright\arraybackslash}p{0.74\linewidth}@{\hspace{0.55em}}>{\raggedright\arraybackslash}p{0.16\linewidth}@{\hspace{0.25em}}}
\mbox{NWGmc Final Long O Raising} & \emph{*bōku} \\
\end{tabular}
\end{minipage}
&
&
\begin{minipage}[t]{\linewidth}
\centering\textbf{Old English changes}\par
\vspace{0.35em}
\raggedright
\begin{tabular}{@{}>{\raggedright\arraybackslash}p{0.74\linewidth}@{\hspace{0.55em}}>{\raggedright\arraybackslash}p{0.16\linewidth}@{\hspace{0.25em}}}
\mbox{OE High Vowel Apocope} & \emph{*bōk} \\
\end{tabular}
\end{minipage}
\\
\end{tabularx}
\end{minipage}%
} \endgroup

Outcome: \emph{bōc}

\paragraph{Reconstruction and comparative
evidence}\label{reconstruction-and-comparative-evidence-1}

Kroonen gives the beech noun as \emph{*bōk(j)ō-} and cites Old English
boc, \emph{bēce} among its reflexes (\citeproc{ref-Kroonen2013}{Kroonen
2013}). The selected input \emph{*bōkō} is the nominative-singular shape
of that family, which is the relevant comparison form here.

\paragraph{Old English evidence}\label{old-english-evidence-1}

Kroonen's Old English evidence already separates the paradigm material:
\emph{boc} as the nominative form and \emph{bēce} as an oblique form
(\citeproc{ref-Kroonen2013}{Kroonen 2013}). The relevant comparator is
therefore \emph{bōc}; \emph{bēċe} remains related paradigm evidence
rather than the form chosen for this comparison.

\paragraph{Development to Old
English}\label{development-to-old-english-1}

With nominative input \emph{*bōkō}, the development is compact.
Northwest Germanic final long \emph{ō} raises to \emph{u}, and later
high-vowel apocope leaves \emph{bōc}. The regular comparison is
therefore \emph{*bōkō} \textgreater{} \emph{bōc}.

\subsubsection{bier --- OE bǣr}\label{bier-oe-bux1e3r}

Derivation: \emph{*bḗrō} \textgreater{} \emph{bǣr} (regular).

\paragraph{Derivation trace}\label{derivation-trace-2}

Proto input: \emph{*bḗrō}

\begingroup
\setlength{\fboxsep}{6pt}

\noindent\fbox{%
\begin{minipage}{0.97\linewidth}
\small
\begin{tabularx}{\linewidth}{@{}>{\raggedright\arraybackslash}p{0.460\linewidth}>{\centering\arraybackslash}X>{\raggedright\arraybackslash}p{0.460\linewidth}@{}}
\begin{minipage}[t]{\linewidth}
\centering\textbf{Earlier Germanic changes}\par
\vspace{0.35em}
\raggedright
\centering\textbf{West Germanic}\par
\raggedright
\vspace{0.2em}
\raggedright [no change]\par
\vspace{0.6em}
\centering\textbf{Northwest Germanic}\par
\raggedright
\vspace{0.2em}
\begin{tabular}{@{}>{\raggedright\arraybackslash}p{0.74\linewidth}@{\hspace{0.55em}}>{\raggedright\arraybackslash}p{0.16\linewidth}@{\hspace{0.25em}}}
\mbox{NWGmc Final Long O Raising} & \emph{*bḗru} \\
\mbox{NWGmc Long E Lowering} & \emph{*bǣru} \\
\end{tabular}
\end{minipage}
&
&
\begin{minipage}[t]{\linewidth}
\centering\textbf{Old English changes}\par
\vspace{0.35em}
\raggedright
\begin{tabular}{@{}>{\raggedright\arraybackslash}p{0.74\linewidth}@{\hspace{0.55em}}>{\raggedright\arraybackslash}p{0.16\linewidth}@{\hspace{0.25em}}}
\mbox{OE High Vowel Apocope} & \emph{*bǣr} \\
\end{tabular}
\end{minipage}
\\
\end{tabularx}
\end{minipage}%
} \endgroup

Outcome: \emph{bǣr}

\paragraph{Reconstruction and comparative
evidence}\label{reconstruction-and-comparative-evidence-2}

Kroonen reconstructs the noun as \emph{*bērō-} f.~`bier' and cites Old
English bar, \emph{bær} among the reflexes
(\citeproc{ref-Kroonen2013}{Kroonen 2013, 717}). The selected input
\emph{*bḗrō} is the same lexeme in the accent notation used here.

\paragraph{Old English evidence}\label{old-english-evidence-2}

Clark Hall and Bosworth-Toller lemmatize the noun as \emph{bær}, and
Kroonen also records \emph{bar} beside it
(\citeproc{ref-ClarkHall1960}{Clark Hall 1960};
\citeproc{ref-BosworthToller1898}{Bosworth and Toller 1898, 73};
\citeproc{ref-Kroonen2013}{Kroonen 2013, 717}). The target \emph{bǣr} is
therefore a normalized long-vowel spelling of the same noun.

\paragraph{Source note}\label{source-note}

Lexicographic spellings vary between \emph{bær} and \emph{bar}. The
normalized target \emph{bǣr} simply marks the same long vowel explicitly
(\citeproc{ref-ClarkHall1960}{Clark Hall 1960};
\citeproc{ref-BosworthToller1898}{Bosworth and Toller 1898, 73};
\citeproc{ref-Kroonen2013}{Kroonen 2013, 717}).

\paragraph{Development to Old
English}\label{development-to-old-english-2}

From \emph{*bḗrō}, Northwest Germanic final long \emph{ō} raises to
\emph{u}, long \emph{ē} lowers to \emph{ǣ}, and high-vowel apocope
yields \emph{bǣr}. The resulting noun matches the normalized Old English
target.

\subsubsection{both --- OE bū}\label{both-oe-bux16b}

Derivation: \emph{*bō} \textgreater{} \emph{bū} (regular).

\paragraph{Derivation trace}\label{derivation-trace-3}

Proto input: \emph{*bō}

\begingroup
\setlength{\fboxsep}{6pt}

\noindent\fbox{%
\begin{minipage}{0.97\linewidth}
\small
\begin{tabularx}{\linewidth}{@{}>{\raggedright\arraybackslash}p{0.540\linewidth}>{\centering\arraybackslash}X>{\raggedright\arraybackslash}p{0.300\linewidth}@{}}
\begin{minipage}[t]{\linewidth}
\centering\textbf{Earlier Germanic changes}\par
\vspace{0.35em}
\raggedright
\centering\textbf{West Germanic}\par
\raggedright
\vspace{0.2em}
\raggedright [no change]\par
\vspace{0.6em}
\centering\textbf{Northwest Germanic}\par
\raggedright
\vspace{0.2em}
\begin{tabular}{@{}>{\raggedright\arraybackslash}p{0.74\linewidth}@{\hspace{0.55em}}>{\raggedright\arraybackslash}p{0.16\linewidth}@{\hspace{0.25em}}}
NWGmc Stressed Monosyllable O Raising & \emph{*bū} \\
\end{tabular}
\end{minipage}
&
&
\begin{minipage}[t]{\linewidth}
\centering\textbf{Old English changes}\par
\vspace{0.35em}
\raggedright
\raggedright [no change]\par
\end{minipage}
\\
\end{tabularx}
\end{minipage}%
} \endgroup

Outcome: \emph{bū}

\paragraph{Reconstruction and comparative
evidence}\label{reconstruction-and-comparative-evidence-3}

Kroonen treats the Germanic numeral under \emph{*ba-} and gives the
inherited paradigm \emph{*bai}, \emph{*bans}, \emph{*bōz}/\emph{*bōns},
\emph{*bō}, with Old English \emph{bēġen}, \emph{bā}, and neuter
\emph{bū} (\citeproc{ref-Kroonen2013}{Kroonen 2013, 47}). For the
present entry, the relevant inherited form is the unextended neuter dual
\emph{*bō}.

The older explanation of \emph{bēġen} derives it from \emph{*bō-jen-},
and Orel still gives OE \emph{bezen} (\textless{} \emph{*bō-jenō)}
beside ON \emph{báðir}, OFris \emph{bēthe}, \emph{OS be-thia}, and OHG
\emph{bēde}, beide (\citeproc{ref-Orel2003}{Orel 2003, 65}). Fulk
reports that explanation cautiously and notes Seebold's preference for a
\emph{*bō-þ-} analysis instead (\citeproc{ref-Fulk2018}{Fulk 2018, sect.
10.1}). That debate matters for \emph{bēġen} and for the extended forms
behind Modern English \emph{both}, German \emph{beide}, and Dutch
\emph{beide}; it does not displace the inherited neuter \emph{*bō}
\textgreater{} \emph{bū} treated here.

\paragraph{Old English evidence}\label{old-english-evidence-3}

The Old English dual paradigm is well established. Brunner gives
masculine \emph{bēġen}, feminine \emph{bā}, and neuter \emph{bū} beside
\emph{bā}, with compounds such as \emph{bā} \emph{twā}, \emph{bū}
\emph{tū}, and \emph{bām} \emph{twām}
(\citeproc{ref-SieversBrunner1965}{Sievers 1965, sect. 324} Anm. 2).
Campbell and Fulk present the same basic pattern: masculine
\emph{bēġen}, feminine \emph{bā}, neuter \emph{bā}, \emph{bū}, genitive
\emph{bēġra}, \emph{bēġ(e)a}, and dative \emph{bǣm}
(\citeproc{ref-Campbell1959}{Campbell 1959, sect. 683};
\citeproc{ref-Fulk2018}{Fulk 2018, sect. 10.1}).

\emph{bū} is therefore an attested neuter dual form, not a
reconstruction. It is the cleanest target for this entry because
\emph{bēġen} belongs to the historically more contested \emph{*bō-jen-}
/ analogical zone, while \emph{bā} remains a partner form within the
dual paradigm rather than the most straightforward monosyllabic
comparison.

\paragraph{Development to Old
English}\label{development-to-old-english-3}

\emph{*bō} is a stressed monosyllabic form. Campbell cites \emph{cū},
\emph{hū}, \emph{tū}, and \emph{bū} as examples of final accented
\emph{ō} \textgreater{} \emph{ū} in the West Germanic stage leading to
Old English (\citeproc{ref-Campbell1959}{Campbell 1959, sect. 122}).
Brunner states the same development more directly: Auslautendes \emph{ō}
erscheint als û in \emph{bū} \ldots{} cu \ldots{} \emph{hū}, \emph{tū}
(\citeproc{ref-SieversBrunner1965}{Sievers 1965, sect. 69}).

The development is therefore straightforward: \emph{*bō} \textgreater{}
\emph{bū}.

\paragraph{Form comparison}\label{form-comparison}

The comparison below is manual. It separates the inherited OE target
from the other forms that belong to the same broader lexical history.

{\def\LTcaptype{none} % do not increment counter
\begin{longtable}[]{@{}
  >{\raggedright\arraybackslash}p{(\linewidth - 6\tabcolsep) * \real{0.2500}}
  >{\raggedright\arraybackslash}p{(\linewidth - 6\tabcolsep) * \real{0.2500}}
  >{\raggedright\arraybackslash}p{(\linewidth - 6\tabcolsep) * \real{0.2500}}
  >{\raggedright\arraybackslash}p{(\linewidth - 6\tabcolsep) * \real{0.2500}}@{}}
\toprule\noalign{}
\begin{minipage}[b]{\linewidth}\raggedright
Form
\end{minipage} & \begin{minipage}[b]{\linewidth}\raggedright
Source / stage
\end{minipage} & \begin{minipage}[b]{\linewidth}\raggedright
Status
\end{minipage} & \begin{minipage}[b]{\linewidth}\raggedright
Relevance to this entry
\end{minipage} \\
\midrule\noalign{}
\endhead
\bottomrule\noalign{}
\endlastfoot
\emph{*bō} \textgreater{} \emph{bū} & PGmc neuter dual \textgreater{} OE
neuter dual & selected regular comparison & main line of the entry \\
\emph{bēġen} & OE masculine dual & attested, but historically contested
and at least partly analogical in Kroonen & real OE evidence, not the
selected target \\
\emph{bā} & OE feminine dual; also neuter variant & attested partner
form & part of the OE paradigm, but not the chosen monosyllabic
comparator \\
\emph{báðir}, \emph{beide}, \emph{both} & Norse, continental West
Germanic, Modern English extended forms & related but different
formation & useful background, not the direct continuation of OE
\emph{bū} \\
\end{longtable}
}

\subsubsection{bow --- OE bīeġan}\label{bow-oe-bux12beux121an}

Derivation: \emph{*báugijaną} \textgreater{} \emph{bīeġan} (regular).

\paragraph{Derivation trace}\label{derivation-trace-4}

Proto input: \emph{*báugijaną}

\begingroup
\setlength{\fboxsep}{6pt}

\noindent\fbox{%
\begin{minipage}{0.97\linewidth}
\footnotesize
\begin{tabularx}{\linewidth}{@{}>{\raggedright\arraybackslash}p{0.280\linewidth}>{\centering\arraybackslash}X>{\raggedright\arraybackslash}p{0.560\linewidth}@{}}
\begin{minipage}[t]{\linewidth}
\centering\textbf{Earlier Germanic changes}\par
\vspace{0.35em}
\raggedright
\centering\textbf{West Germanic}\par
\raggedright
\vspace{0.2em}
\raggedright [no change]\par
\vspace{0.6em}
\centering\textbf{Northwest Germanic}\par
\raggedright
\vspace{0.2em}
\raggedright [no change]\par
\end{minipage}
&
&
\begin{minipage}[t]{\linewidth}
\centering\textbf{Old English changes}\par
\vspace{0.35em}
\raggedright
\begin{tabular}{@{}>{\raggedright\arraybackslash}p{0.62\linewidth}@{\hspace{0.55em}}>{\raggedright\arraybackslash}p{0.28\linewidth}@{\hspace{0.25em}}}
OE Au Fronting & \emph{*báeugijaną} \\
OE Diphthong Leveling & \emph{*bēagijaną} \\
OE Heavy Syllable Nasal Apocope & \emph{*bēagijan} \\
OE Secondary Nasalization & \emph{*bēagijąn} \\
Sievers Law Syncope & \emph{*bēagjąn} \\
OE Velar Palatalization & \emph{*bēaʤjąn} \\
OE I Umlaut & \emph{*bīeʤjąn} \\
OE Weak Tail Reduction & \emph{*bīeʤjan} \\
OE J Loss After Heavy & \emph{*bīeʤan} \\
\end{tabular}
\end{minipage}
\\
\end{tabularx}
\end{minipage}%
} \endgroup

Outcome: \emph{bīeġan}

\paragraph{Reconstruction and comparative
evidence}\label{reconstruction-and-comparative-evidence-4}

Kroonen reconstructs the weak verb as \emph{*baugjan-} `to (make) bend'
and cites Old English \emph{biegan} among its reflexes
(\citeproc{ref-Kroonen2013}{Kroonen 2013}). Ringe and Taylor give the
northwest Germanic and West Saxon development more fully as PNWGmce
\emph{*baugijana} \textgreater{} \emph{*béagjan} \textgreater{} WS OE
\emph{biegan} (\citeproc{ref-RingeTaylor2014}{Ringe and Taylor 2014}).
The entry therefore concerns the weak causative member of the
bend-family, alongside the related strong verb and noun.

\paragraph{Old English evidence}\label{old-english-evidence-4}

Clark Hall lemmatizes \emph{biegan}, and Bosworth-Toller records
\emph{bigan} with examples such as Ic \emph{bēge} \emph{mīne}
\emph{cneówa} and Se ord \emph{bīgde} upp \emph{tō} \emph{þām} hiltum
(\citeproc{ref-ClarkHall1960}{Clark Hall 1960};
\citeproc{ref-BosworthToller1898}{Bosworth and Toller 1898, 102}). The
form \emph{bīeġan} used here is a normalized spelling of that attested
Old English weak verb.

\paragraph{Development to Old
English}\label{development-to-old-english-4}

From \emph{*báugijaną}, the stem reaches pre-Old-English
\emph{*bēagjan}, after which palatalization of \emph{*gj} and i-umlaut
yield West Saxon \emph{biegan}; Campbell lists \emph{biegan} among the
regular \emph{ie} outcomes of \emph{*éa} under i-umlaut
(\citeproc{ref-RingeTaylor2014}{Ringe and Taylor 2014};
\citeproc{ref-Campbell1959}{Campbell 1959, 80}). The development is
therefore straightforward: \emph{*báugijaną} \textgreater{}
\emph{bīeġan}.

\subsubsection{fare --- OE faran}\label{fare-oe-faran}

Derivation: \emph{*fáraną} \textgreater{} \emph{faran} (regular).

\paragraph{Derivation trace}\label{derivation-trace-5}

Proto input: \emph{*fáraną}

\begingroup
\setlength{\fboxsep}{6pt}

\noindent\fbox{%
\begin{minipage}{0.97\linewidth}
\small
\begin{tabularx}{\linewidth}{@{}>{\raggedright\arraybackslash}p{0.300\linewidth}>{\centering\arraybackslash}X>{\raggedright\arraybackslash}p{0.540\linewidth}@{}}
\begin{minipage}[t]{\linewidth}
\centering\textbf{Earlier Germanic changes}\par
\vspace{0.35em}
\raggedright
\centering\textbf{West Germanic}\par
\raggedright
\vspace{0.2em}
\raggedright [no change]\par
\vspace{0.6em}
\centering\textbf{Northwest Germanic}\par
\raggedright
\vspace{0.2em}
\raggedright [no change]\par
\end{minipage}
&
&
\begin{minipage}[t]{\linewidth}
\centering\textbf{Old English changes}\par
\vspace{0.35em}
\raggedright
\begin{tabular}{@{}>{\raggedright\arraybackslash}p{0.74\linewidth}@{\hspace{0.55em}}>{\raggedright\arraybackslash}p{0.16\linewidth}@{\hspace{0.25em}}}
Anglo Frisian Brightening & \emph{*færaną} \\
OE A Restoration & \emph{*faraną} \\
OE Heavy Syllable Nasal Apocope & \emph{*faran} \\
OE Secondary Nasalization & \emph{*farąn} \\
OE Weak Tail Reduction & \emph{*faran} \\
\end{tabular}
\end{minipage}
\\
\end{tabularx}
\end{minipage}%
} \endgroup

Outcome: \emph{faran}

\paragraph{Reconstruction and comparative
evidence}\label{reconstruction-and-comparative-evidence-5}

Kroonen gives the inherited strong verb as \emph{*faran-}, and Orel
gives the same lexeme as \emph{*faranan}, both with Old English
\emph{faran} among the reflexes (\citeproc{ref-Kroonen2013}{Kroonen
2013}; \citeproc{ref-Orel2003}{Orel 2003, 132}). Campbell also uses
\emph{faran} as a standard example of Old English A-restoration
(\citeproc{ref-Campbell1959}{Campbell 1959, 61}).

\paragraph{Old English evidence}\label{old-english-evidence-5}

Clark Hall lemmatizes the strong verb as \emph{faran} and separately
records weak \emph{færan} `to frighten'; \emph{fære}, \emph{færst}, and
\emph{færð} belong to present-tense forms of \emph{faran} rather than to
the infinitive itself (\citeproc{ref-ClarkHall1960}{Clark Hall 1960}).
Bosworth-Toller preserves the same distinction
(\citeproc{ref-BosworthToller1898}{Bosworth and Toller 1898, 108}). The
selected target is therefore the attested citation infinitive
\emph{faran}.

\paragraph{Development to Old
English}\label{development-to-old-english-5}

From \emph{*fáraną}, Anglo-Frisian brightening first gives
\emph{*færaną}, but A-restoration before single \emph{r} returns
\emph{*faraną}; later apocope and weak-tail reduction yield \emph{faran}
(\citeproc{ref-Campbell1959}{Campbell 1959, 61}). Fulk's contrast with
participial faren- \textless{} \emph{*faræn-} \textless{} \emph{*faran-}
shows why fronting elsewhere in the paradigm does not alter the
infinitive headword (\citeproc{ref-Fulk2018}{Fulk 2018}).

\subsubsection{guest --- OE ġiest}\label{guest-oe-ux121iest}

Derivation: \emph{*gástiz} \textgreater{} \emph{ġiest} (regular).

\paragraph{Derivation trace}\label{derivation-trace-6}

Proto input: \emph{*gástiz}

\begingroup
\setlength{\fboxsep}{6pt}

\noindent\fbox{%
\begin{minipage}{0.97\linewidth}
\small
\begin{tabularx}{\linewidth}{@{}>{\raggedright\arraybackslash}p{0.460\linewidth}>{\centering\arraybackslash}X>{\raggedright\arraybackslash}p{0.460\linewidth}@{}}
\begin{minipage}[t]{\linewidth}
\centering\textbf{Earlier Germanic changes}\par
\vspace{0.35em}
\raggedright
\centering\textbf{West Germanic}\par
\raggedright
\vspace{0.2em}
\raggedright [no change]\par
\vspace{0.6em}
\centering\textbf{Northwest Germanic}\par
\raggedright
\vspace{0.2em}
\begin{tabular}{@{}>{\raggedright\arraybackslash}p{0.74\linewidth}@{\hspace{0.55em}}>{\raggedright\arraybackslash}p{0.16\linewidth}@{\hspace{0.25em}}}
\mbox{PGmc Final Z Deletion} & \emph{*gásti} \\
\end{tabular}
\end{minipage}
&
&
\begin{minipage}[t]{\linewidth}
\centering\textbf{Old English changes}\par
\vspace{0.35em}
\raggedright
\begin{tabular}{@{}>{\raggedright\arraybackslash}p{0.74\linewidth}@{\hspace{0.55em}}>{\raggedright\arraybackslash}p{0.16\linewidth}@{\hspace{0.25em}}}
Anglo Frisian Brightening & \emph{*gæsti} \\
OE Velar Palatalization & \emph{*ʤæsti} \\
OE I Umlaut & \emph{*ʤesti} \\
OE Ws Palatal Diphthongization & \emph{*ʤiesti} \\
\mbox{OE High Vowel Apocope} & \emph{*ʤiest} \\
\end{tabular}
\end{minipage}
\\
\end{tabularx}
\end{minipage}%
} \endgroup

Outcome: \emph{ġiest}

\paragraph{Reconstruction and comparative
evidence}\label{reconstruction-and-comparative-evidence-6}

Campbell and Ringe-Taylor treat the noun as an ordinary i-stem whose
West Saxon development shows palatal diphthongization, while
non-West-Saxon evidence preserves forms of the \emph{gest} type
(\citeproc{ref-Campbell1959}{Campbell 1959};
\citeproc{ref-RingeTaylor2014}{Ringe and Taylor 2014}).

\paragraph{Old English evidence}\label{old-english-evidence-6}

Bosworth-Toller and Clark Hall record the word under forms such as
\emph{gist}, \emph{gest}, \emph{giest}, and \emph{gyst}
(\citeproc{ref-BosworthToller1898}{Bosworth and Toller 1898};
\citeproc{ref-ClarkHall1960}{Clark Hall 1960}). The selected target
\emph{ġiest} is the normalized West Saxon form within that attested
family.

\paragraph{Development to Old
English}\label{development-to-old-english-6}

From \emph{*gástiz}, Anglo-Frisian brightening gives a \emph{gæst-}
stage, and i-mutation affects the front vowel before the lost
high-vocalic ending. In West Saxon the initial palatal environment then
produces \emph{ie}, so the regular outcome is \emph{ġiest}
(\citeproc{ref-Campbell1959}{Campbell 1959};
\citeproc{ref-RingeTaylor2014}{Ringe and Taylor 2014}).

\paragraph{Dialect note}\label{dialect-note}

West Saxon \emph{ġiest} is the selected target here. Anglian \emph{gest}
and related spellings remain real Old English comparators rather than
corrections to that choice (\citeproc{ref-RingeTaylor2014}{Ringe and
Taylor 2014}; \citeproc{ref-BosworthToller1898}{Bosworth and Toller
1898}).

\subsubsection{learn --- OE liornian}\label{learn-oe-liornian}

Derivation: \emph{*líznōjaną} \textgreater{} \emph{liornian} (regular).

\paragraph{Derivation trace}\label{derivation-trace-7}

Proto input: \emph{*líznōjaną}

\begingroup
\setlength{\fboxsep}{6pt}

\noindent\fbox{%
\begin{minipage}{0.97\linewidth}
\footnotesize
\begin{tabularx}{\linewidth}{@{}>{\raggedright\arraybackslash}p{0.280\linewidth}>{\centering\arraybackslash}X>{\raggedright\arraybackslash}p{0.560\linewidth}@{}}
\begin{minipage}[t]{\linewidth}
\centering\textbf{Earlier Germanic changes}\par
\vspace{0.35em}
\raggedright
\centering\textbf{West Germanic}\par
\raggedright
\vspace{0.2em}
\raggedright [no change]\par
\vspace{0.6em}
\centering\textbf{Northwest Germanic}\par
\raggedright
\vspace{0.2em}
\raggedright [no change]\par
\end{minipage}
&
&
\begin{minipage}[t]{\linewidth}
\centering\textbf{Old English changes}\par
\vspace{0.35em}
\raggedright
\begin{tabular}{@{}>{\raggedright\arraybackslash}p{0.62\linewidth}@{\hspace{0.55em}}>{\raggedright\arraybackslash}p{0.28\linewidth}@{\hspace{0.25em}}}
OE Breaking & \emph{*líornōjaną} \\
OE Heavy Syllable Nasal Apocope & \emph{*líornōjan} \\
OE Secondary Nasalization & \emph{*líornōjąn} \\
OE I Umlaut & \emph{*líornējąn} \\
OE Unstressed Long Vowel Shortening & \emph{*líornejąn} \\
OE Weak Tail Reduction & \emph{*líornejan} \\
OE Intervocalic J Vocalization & \emph{*líorneian} \\
OE Unstressed EI Contraction & \emph{*líornian} \\
\end{tabular}
\end{minipage}
\\
\end{tabularx}
\end{minipage}%
} \endgroup

Outcome: \emph{liornian}

\paragraph{Reconstruction and comparative
evidence}\label{reconstruction-and-comparative-evidence-7}

Ringe and Taylor place the verb in a class-II weak family of the
\emph{*liznō-} type (\citeproc{ref-RingeTaylor2014}{Ringe and Taylor
2014}), and Kroonen keeps the same comparative base for the Germanic
lexeme (\citeproc{ref-Kroonen2013}{Kroonen 2013}). The selected input
therefore requires no change of stem class or paradigm cell.

The non-obvious issue is dialectal. Campbell records Northumbrian
\emph{liornian} beside \emph{leornian}, and he states that the \emph{eo}
of \emph{leornian} is secondary, while Northumbrian preserves forms with
\emph{io}, where original \emph{eo} and \emph{io} remain distinct
(\citeproc{ref-Campbell1959}{Campbell 1959, sect. 123} n.~2).

\paragraph{Old English evidence}\label{old-english-evidence-7}

The form modeled here is \emph{liornian}, the Northumbrian member of the
Old English family. Dictionary practice more often privileges
\emph{leornian} as the ordinary headword
(\citeproc{ref-ClarkHall1960}{Clark Hall 1960};
\citeproc{ref-BrightCassidyRingler1971}{Bright 1971}), but Campbell's
dialect evidence shows that \emph{liornian} is a genuine Old English
form (\citeproc{ref-Campbell1959}{Campbell 1959, sect. 123} n.~2).

This entry therefore remains compact. The point is to state clearly that
the selected target belongs to the Northumbrian side of the OE evidence
rather than to the leveled \emph{leornian} headword tradition.

\paragraph{Development to Old
English}\label{development-to-old-english-7}

From \emph{*líznōjaną}, the expected Old English developments include
rhotacism of \emph{z}, followed by breaking before \emph{r} plus
consonant, and the ordinary reduction of the weak verbal ending. The
result is \emph{liornian}, preserving the \emph{io} spelling that
Campbell associates with Northumbrian where original \emph{eo} and
\emph{io} remain distinct (\citeproc{ref-Campbell1959}{Campbell 1959,
sect. 123} n.~2).

\emph{Leornian} reflects the later \emph{eo} development that Campbell
treats as secondary {[}Campbell (\citeproc{ref-Campbell1959}{1959}),
§123 n.~2; §296{]}. The selected Northumbrian target is therefore the
regular comparison form for the \emph{i}-grade member of the family.

\paragraph{Form comparison}\label{form-comparison-1}

The comparison below is manual. It distinguishes the regular
Northumbrian form modeled here from the better-known West Saxon
headword.

{\def\LTcaptype{none} % do not increment counter
\begin{longtable}[]{@{}
  >{\raggedright\arraybackslash}p{(\linewidth - 4\tabcolsep) * \real{0.3333}}
  >{\raggedright\arraybackslash}p{(\linewidth - 4\tabcolsep) * \real{0.3333}}
  >{\raggedright\arraybackslash}p{(\linewidth - 4\tabcolsep) * \real{0.3333}}@{}}
\toprule\noalign{}
\begin{minipage}[b]{\linewidth}\raggedright
Form
\end{minipage} & \begin{minipage}[b]{\linewidth}\raggedright
Status
\end{minipage} & \begin{minipage}[b]{\linewidth}\raggedright
Relevance to this entry
\end{minipage} \\
\midrule\noalign{}
\endhead
\bottomrule\noalign{}
\endlastfoot
\emph{*líznōjaną} -\textgreater{} liornian & computed regular output;
attested Northumbrian comparison form & selected comparison \\
\emph{leornian} & attested later \emph{eo} form and dictionary headword
& useful control, but not the target of this entry \\
\end{longtable}
}

\subsubsection{linden --- OE lind}\label{linden-oe-lind}

Derivation: \emph{*líndō} \textgreater{} \emph{lind} (regular).

\paragraph{Derivation trace}\label{derivation-trace-8}

Proto input: \emph{*líndō}

\begingroup
\setlength{\fboxsep}{6pt}

\noindent\fbox{%
\begin{minipage}{0.97\linewidth}
\small
\begin{tabularx}{\linewidth}{@{}>{\raggedright\arraybackslash}p{0.460\linewidth}>{\centering\arraybackslash}X>{\raggedright\arraybackslash}p{0.460\linewidth}@{}}
\begin{minipage}[t]{\linewidth}
\centering\textbf{Earlier Germanic changes}\par
\vspace{0.35em}
\raggedright
\centering\textbf{West Germanic}\par
\raggedright
\vspace{0.2em}
\raggedright [no change]\par
\vspace{0.6em}
\centering\textbf{Northwest Germanic}\par
\raggedright
\vspace{0.2em}
\begin{tabular}{@{}>{\raggedright\arraybackslash}p{0.74\linewidth}@{\hspace{0.55em}}>{\raggedright\arraybackslash}p{0.16\linewidth}@{\hspace{0.25em}}}
\mbox{NWGmc Final Long O Raising} & \emph{*líndu} \\
\end{tabular}
\end{minipage}
&
&
\begin{minipage}[t]{\linewidth}
\centering\textbf{Old English changes}\par
\vspace{0.35em}
\raggedright
\begin{tabular}{@{}>{\raggedright\arraybackslash}p{0.74\linewidth}@{\hspace{0.55em}}>{\raggedright\arraybackslash}p{0.16\linewidth}@{\hspace{0.25em}}}
\mbox{OE High Vowel Apocope} & \emph{*línd} \\
\end{tabular}
\end{minipage}
\\
\end{tabularx}
\end{minipage}%
} \endgroup

Outcome: \emph{lind}

\paragraph{Reconstruction and comparative
evidence}\label{reconstruction-and-comparative-evidence-8}

Kroonen cites \emph{*lindō-} `lime tree' and gives Old English
\emph{lind} as the relevant reflex (\citeproc{ref-Kroonen2013}{Kroonen
2013}).

\paragraph{Old English evidence}\label{old-english-evidence-8}

Clark Hall and Bosworth-Toller both record \emph{lind} as the noun
`lime-tree, linden' (\citeproc{ref-ClarkHall1960}{Clark Hall 1960};
\citeproc{ref-BosworthToller1898}{Bosworth and Toller 1898, 630}).

\paragraph{Development to Old
English}\label{development-to-old-english-8}

From \emph{*líndō}, Northwest Germanic final \emph{*ō} raising first
gives a \emph{*líndu} stage, and later high-vowel apocope yields
\emph{lind}. The development is therefore straightforward and regular.

\paragraph{Form note}\label{form-note}

The Old English noun represented here is \emph{lind}. Clark Hall also
has a separate adjectival \emph{linden} `made of linden-wood', but that
is not the noun counterpart for this entry
(\citeproc{ref-ClarkHall1960}{Clark Hall 1960}).

\subsubsection{mother --- OE mōder}\label{mother-oe-mux14dder}

Derivation: \emph{*mōdēr} \textgreater{} \emph{mōder} (regular).

\paragraph{Derivation trace}\label{derivation-trace-9}

Proto input: \emph{*mōdēr}

\begingroup
\setlength{\fboxsep}{6pt}

\noindent\fbox{%
\begin{minipage}{0.97\linewidth}
\small
\begin{tabularx}{\linewidth}{@{}>{\raggedright\arraybackslash}p{0.460\linewidth}>{\centering\arraybackslash}X>{\raggedright\arraybackslash}p{0.460\linewidth}@{}}
\begin{minipage}[t]{\linewidth}
\centering\textbf{Earlier Germanic changes}\par
\vspace{0.35em}
\raggedright
\centering\textbf{West Germanic}\par
\raggedright
\vspace{0.2em}
\raggedright [no change]\par
\vspace{0.6em}
\centering\textbf{Northwest Germanic}\par
\raggedright
\vspace{0.2em}
\begin{tabular}{@{}>{\raggedright\arraybackslash}p{0.74\linewidth}@{\hspace{0.55em}}>{\raggedright\arraybackslash}p{0.16\linewidth}@{\hspace{0.25em}}}
\mbox{NWGmc Long E Lowering} & \emph{*mōdǣr} \\
\end{tabular}
\end{minipage}
&
&
\begin{minipage}[t]{\linewidth}
\centering\textbf{Old English changes}\par
\vspace{0.35em}
\raggedright
\begin{tabular}{@{}>{\raggedright\arraybackslash}p{0.74\linewidth}@{\hspace{0.55em}}>{\raggedright\arraybackslash}p{0.16\linewidth}@{\hspace{0.25em}}}
OE Unstressed Long Vowel Shortening & \emph{*mōdær} \\
OE Unstressed AE Merger & \emph{*mōder} \\
\end{tabular}
\end{minipage}
\\
\end{tabularx}
\end{minipage}%
} \endgroup

Outcome: \emph{mōder}

\paragraph{Reconstruction and comparative
evidence}\label{reconstruction-and-comparative-evidence-9}

Kroonen and Orel cite the Proto-Germanic r-stem kinship noun as
\emph{*mōder-} / \emph{*mōdēr} (\citeproc{ref-Kroonen2013}{Kroonen
2013}; \citeproc{ref-Orel2003}{Orel 2003}).

\paragraph{Old English evidence}\label{old-english-evidence-9}

The transmitted Old English headword tradition is \emph{mōdor} / modor,
with oblique \emph{mēder} in the paradigm. Clark Hall, Campbell, and
Ringe and Taylor all preserve that contrast
(\citeproc{ref-ClarkHall1960}{Clark Hall 1960};
\citeproc{ref-Campbell1959}{Campbell 1959};
\citeproc{ref-RingeTaylor2014}{Ringe and Taylor 2014}).

\paragraph{Development to Old
English}\label{development-to-old-english-9}

From \emph{*mōdēr}, the regular suffixal development yields
\emph{mōder}. That regular nominative reflex is the form represented
here, while the more familiar citation form \emph{mōdor} reflects later
levelling within the r-stem paradigm
(\citeproc{ref-Campbell1959}{Campbell 1959};
\citeproc{ref-RingeTaylor2014}{Ringe and Taylor 2014}).

\paragraph{Form comparison}\label{form-comparison-2}

The note therefore concerns inherited vocalism rather than a different
lexeme: \emph{mōder} is the regularized nominative represented here, but
dictionaries usually print \emph{mōdor} / modor, and the oblique
evidence survives in \emph{mēder} (\citeproc{ref-ClarkHall1960}{Clark
Hall 1960}; \citeproc{ref-RingeTaylor2014}{Ringe and Taylor 2014}).

\subsubsection{show --- OE sċēawian}\label{show-oe-sux10bux113awian}

Derivation: \emph{*skáwōjaną} \textgreater{} \emph{sċēawian} (regular).

\paragraph{Derivation trace}\label{derivation-trace-10}

Proto input: \emph{*skáwōjaną}

\begingroup
\setlength{\fboxsep}{6pt}

\noindent\fbox{%
\begin{minipage}{0.97\linewidth}
\footnotesize
\begin{tabularx}{\linewidth}{@{}>{\raggedright\arraybackslash}p{0.280\linewidth}>{\centering\arraybackslash}X>{\raggedright\arraybackslash}p{0.560\linewidth}@{}}
\begin{minipage}[t]{\linewidth}
\centering\textbf{Earlier Germanic changes}\par
\vspace{0.35em}
\raggedright
\centering\textbf{West Germanic}\par
\raggedright
\vspace{0.2em}
\raggedright [no change]\par
\vspace{0.6em}
\centering\textbf{Northwest Germanic}\par
\raggedright
\vspace{0.2em}
\raggedright [no change]\par
\end{minipage}
&
&
\begin{minipage}[t]{\linewidth}
\centering\textbf{Old English changes}\par
\vspace{0.35em}
\raggedright
\begin{tabular}{@{}>{\raggedright\arraybackslash}p{0.62\linewidth}@{\hspace{0.55em}}>{\raggedright\arraybackslash}p{0.28\linewidth}@{\hspace{0.25em}}}
OE Aw Long Diphthong & \emph{*skḗawōjaną} \\
OE Heavy Syllable Nasal Apocope & \emph{*skḗawōjan} \\
OE Secondary Nasalization & \emph{*skḗawōjąn} \\
OE Sk Palatalization & \emph{*ʃḗawōjąn} \\
OE I Umlaut & \emph{*ʃḗawējąn} \\
OE Unstressed Long Vowel Shortening & \emph{*ʃḗawejąn} \\
OE Weak Tail Reduction & \emph{*ʃḗawejan} \\
OE Intervocalic J Vocalization & \emph{*ʃḗaweian} \\
OE Unstressed EI Contraction & \emph{*ʃḗawian} \\
\end{tabular}
\end{minipage}
\\
\end{tabularx}
\end{minipage}%
} \endgroup

Outcome: \emph{sċēawian}

\paragraph{Reconstruction and comparative
evidence}\label{reconstruction-and-comparative-evidence-10}

Orel and Kroonen cite a Class II verb of the type \emph{*skawōjan-},
with OE \emph{scēawian} among the reflexes (\citeproc{ref-Orel2003}{Orel
2003}; \citeproc{ref-Kroonen2013}{Kroonen 2013, 482}). Brunner likewise
records the Old English family as \emph{scēawian}, \emph{scāwian}, which
places this entry in the ordinary show-verb set rather than in a special
finite-cell workaround (\citeproc{ref-SieversBrunner1965}{Sievers
1965}).

\paragraph{Old English evidence}\label{old-english-evidence-10}

Bright lists \emph{scēawian} (W. II.) and also the related form
\emph{scēawa} (\citeproc{ref-BrightCassidyRingler1971}{Bright 1971}).
The source tradition therefore uses \emph{scēawian}, while the target
represented here is the normalized project spelling \emph{sċēawian}.

\paragraph{Development to Old
English}\label{development-to-old-english-10}

From \emph{*skáwōjaną}, Old English \emph{aw} before a following vowel
yields \emph{ēaw}, and the Class II suffix keeps \emph{*ō} between
\emph{*w} and \emph{*j}. The development therefore runs regularly to
\emph{sċēawian}, without the direct \emph{*aw+j} problem seen in other
verb types (\citeproc{ref-Campbell1959}{Campbell 1959};
\citeproc{ref-Orel2003}{Orel 2003}).

\paragraph{Form note}\label{form-note-1}

The difference between \emph{scēawian} and \emph{sċēawian} is
orthographic normalization of initial \emph{}, not a difference of
lexeme or paradigm cell (\citeproc{ref-Campbell1959}{Campbell 1959};
\citeproc{ref-Hogg1992}{Hogg 1992}).

\subsubsection{weapon --- OE wǣpn}\label{weapon-oe-wux1e3pn}

Derivation: \emph{*wḗpną} \textgreater{} \emph{wǣpn} (regular).

\paragraph{Derivation trace}\label{derivation-trace-11}

Proto input: \emph{*wḗpną}

\begingroup
\setlength{\fboxsep}{6pt}

\noindent\fbox{%
\begin{minipage}{0.97\linewidth}
\small
\begin{tabularx}{\linewidth}{@{}>{\raggedright\arraybackslash}p{0.460\linewidth}>{\centering\arraybackslash}X>{\raggedright\arraybackslash}p{0.460\linewidth}@{}}
\begin{minipage}[t]{\linewidth}
\centering\textbf{Earlier Germanic changes}\par
\vspace{0.35em}
\raggedright
\centering\textbf{West Germanic}\par
\raggedright
\vspace{0.2em}
\raggedright [no change]\par
\vspace{0.6em}
\centering\textbf{Northwest Germanic}\par
\raggedright
\vspace{0.2em}
\begin{tabular}{@{}>{\raggedright\arraybackslash}p{0.74\linewidth}@{\hspace{0.55em}}>{\raggedright\arraybackslash}p{0.16\linewidth}@{\hspace{0.25em}}}
\mbox{NWGmc Long E Lowering} & \emph{*wǣpną} \\
\end{tabular}
\end{minipage}
&
&
\begin{minipage}[t]{\linewidth}
\centering\textbf{Old English changes}\par
\vspace{0.35em}
\raggedright
\begin{tabular}{@{}>{\raggedright\arraybackslash}p{0.74\linewidth}@{\hspace{0.55em}}>{\raggedright\arraybackslash}p{0.16\linewidth}@{\hspace{0.25em}}}
OE Heavy Syllable Nasal Apocope & \emph{*wǣpn} \\
\end{tabular}
\end{minipage}
\\
\end{tabularx}
\end{minipage}%
} \endgroup

Outcome: \emph{wǣpn}

\paragraph{Reconstruction and comparative
evidence}\label{reconstruction-and-comparative-evidence-11}

Kroonen reconstructs a double-stem noun \emph{*wēbna-} \textasciitilde{}
\emph{*wēpna-} and cites OE \emph{wæpn} among its reflexes
(\citeproc{ref-Kroonen2013}{Kroonen 2013, 617}). The selected input
\emph{*wḗpną} represents the unbroken citation-form noun rather than the
later broken simplex.

\paragraph{Old English evidence}\label{old-english-evidence-11}

Campbell's cluster-noun discussion preserves unbroken \emph{wépn} beside
broken \emph{wépen}-type forms (\citeproc{ref-Campbell1959}{Campbell
1959, 150}; \citeproc{ref-Campbell1959}{1959, 226--27}). Bright
contrasts broken nominative \emph{wǣpen}/wapen with unbroken oblique
\emph{wǣpnes}, while Clark Hall lemmatizes the noun under \emph{wapen}
and also preserves unbroken forms in compounds and related spellings
(\citeproc{ref-BrightCassidyRingler1971}{Bright 1971, 29};
\citeproc{ref-ClarkHall1960}{Clark Hall 1960, 355}).

\paragraph{Development to Old
English}\label{development-to-old-english-11}

Northwest Germanic lowering gives \emph{wǣpn}, and loss of the final
nasal vowel leaves the unbroken cluster word-finally. The selected
target is the attested unbroken form \emph{wǣpn}.

\paragraph{Form note}\label{form-note-2}

The ordinary late West Saxon simplex headword is \emph{wǣpen}, but
Campbell's noun-class discussion also preserves unbroken \emph{wépn}
beside broken \emph{wépen}-type forms
(\citeproc{ref-Campbell1959}{Campbell 1959, 150};
\citeproc{ref-Campbell1959}{1959, 226--27}). \emph{wǣpnes} remains the
regular unbroken oblique comparator
(\citeproc{ref-BrightCassidyRingler1971}{Bright 1971, 29};
\citeproc{ref-ClarkHall1960}{Clark Hall 1960, 355}).

\clearpage

\subsection{Variant B: compact regular
style}\label{variant-b-compact-regular-style}

These entries keep the heading, derivation line, and boxed trace, but
merge the three standard prose sections into one commentary block and
retain only shorter note/comparison sections separately.

\subsubsection{bake --- OE bacan}\label{bake-oe-bacan-1}

Derivation: \emph{*bákaną} \textgreater{} \emph{bacan} (regular).

\paragraph{Derivation trace}\label{derivation-trace-12}

Proto input: \emph{*bákaną}

\begingroup
\setlength{\fboxsep}{6pt}

\noindent\fbox{%
\begin{minipage}{0.97\linewidth}
\small
\begin{tabularx}{\linewidth}{@{}>{\raggedright\arraybackslash}p{0.300\linewidth}>{\centering\arraybackslash}X>{\raggedright\arraybackslash}p{0.540\linewidth}@{}}
\begin{minipage}[t]{\linewidth}
\centering\textbf{Earlier Germanic changes}\par
\vspace{0.35em}
\raggedright
\centering\textbf{West Germanic}\par
\raggedright
\vspace{0.2em}
\raggedright [no change]\par
\vspace{0.6em}
\centering\textbf{Northwest Germanic}\par
\raggedright
\vspace{0.2em}
\raggedright [no change]\par
\end{minipage}
&
&
\begin{minipage}[t]{\linewidth}
\centering\textbf{Old English changes}\par
\vspace{0.35em}
\raggedright
\begin{tabular}{@{}>{\raggedright\arraybackslash}p{0.74\linewidth}@{\hspace{0.55em}}>{\raggedright\arraybackslash}p{0.16\linewidth}@{\hspace{0.25em}}}
Anglo Frisian Brightening & \emph{*bækaną} \\
OE A Restoration & \emph{*bakaną} \\
OE Heavy Syllable Nasal Apocope & \emph{*bakan} \\
OE Secondary Nasalization & \emph{*bakąn} \\
OE Weak Tail Reduction & \emph{*bakan} \\
\end{tabular}
\end{minipage}
\\
\end{tabularx}
\end{minipage}%
} \endgroup

\paragraph{Commentary}\label{commentary}

Orel reconstructs the verb as \emph{*bakanan} and cites Old English
\emph{bacan} beside Old High German \emph{backan, bahhan}
(\citeproc{ref-Orel2003}{Orel 2003}). Campbell gives \emph{bacan} as one
of the standard examples of Old English A-restoration before a single
consonant, and Ringe and Taylor state the same development from
\emph{*bakan} to Old English \emph{bacan}
(\citeproc{ref-Campbell1959}{Campbell 1959, 61};
\citeproc{ref-RingeTaylor2014}{Ringe and Taylor 2014}).

Bosworth-Toller and Clark Hall both record \emph{bacan} as the ordinary
Old English verb `to bake' (\citeproc{ref-BosworthToller1898}{Bosworth
and Toller 1898, 72}; \citeproc{ref-ClarkHall1960}{Clark Hall 1960}).
The target in this entry is therefore the attested infinitive headword
itself, not a selected oblique or finite paradigm cell.

From \emph{*bákaną}, Anglo-Frisian brightening first gives
\emph{*bækaną}. A-restoration then returns the stem vowel to \emph{a}
before single \emph{k} plus the back-vocalic infinitive suffix, and
later apocope and weak-tail reduction yield \emph{bacan}
(\citeproc{ref-Campbell1959}{Campbell 1959, 61};
\citeproc{ref-RingeTaylor2014}{Ringe and Taylor 2014}). The development
is therefore straightforward: \emph{*bákaną} \textgreater{} bacan.

\subsubsection{beech --- OE bōc}\label{beech-oe-bux14dc-1}

Derivation: \emph{*bōkō} \textgreater{} \emph{bōc} (regular).

\paragraph{Derivation trace}\label{derivation-trace-13}

Proto input: \emph{*bōkō}

\begingroup
\setlength{\fboxsep}{6pt}

\noindent\fbox{%
\begin{minipage}{0.97\linewidth}
\small
\begin{tabularx}{\linewidth}{@{}>{\raggedright\arraybackslash}p{0.460\linewidth}>{\centering\arraybackslash}X>{\raggedright\arraybackslash}p{0.460\linewidth}@{}}
\begin{minipage}[t]{\linewidth}
\centering\textbf{Earlier Germanic changes}\par
\vspace{0.35em}
\raggedright
\centering\textbf{West Germanic}\par
\raggedright
\vspace{0.2em}
\raggedright [no change]\par
\vspace{0.6em}
\centering\textbf{Northwest Germanic}\par
\raggedright
\vspace{0.2em}
\begin{tabular}{@{}>{\raggedright\arraybackslash}p{0.74\linewidth}@{\hspace{0.55em}}>{\raggedright\arraybackslash}p{0.16\linewidth}@{\hspace{0.25em}}}
\mbox{NWGmc Final Long O Raising} & \emph{*bōku} \\
\end{tabular}
\end{minipage}
&
&
\begin{minipage}[t]{\linewidth}
\centering\textbf{Old English changes}\par
\vspace{0.35em}
\raggedright
\begin{tabular}{@{}>{\raggedright\arraybackslash}p{0.74\linewidth}@{\hspace{0.55em}}>{\raggedright\arraybackslash}p{0.16\linewidth}@{\hspace{0.25em}}}
\mbox{OE High Vowel Apocope} & \emph{*bōk} \\
\end{tabular}
\end{minipage}
\\
\end{tabularx}
\end{minipage}%
} \endgroup

\paragraph{Commentary}\label{commentary-1}

Kroonen gives the beech noun as \emph{*bōk(j)ō-} and cites Old English
boc, \emph{bēce} among its reflexes (\citeproc{ref-Kroonen2013}{Kroonen
2013}). The selected input \emph{*bōkō} is the nominative-singular shape
of that family, which is the relevant comparison form here.

Kroonen's Old English evidence already separates the paradigm material:
\emph{boc} as the nominative form and \emph{bēce} as an oblique form
(\citeproc{ref-Kroonen2013}{Kroonen 2013}). The relevant comparator is
therefore \emph{bōc}; \emph{bēċe} remains related paradigm evidence
rather than the form chosen for this comparison.

With nominative input \emph{*bōkō}, the development is compact.
Northwest Germanic final long \emph{ō} raises to \emph{u}, and later
high-vowel apocope leaves \emph{bōc}. The regular comparison is
therefore \emph{*bōkō} \textgreater{} \emph{bōc}.

\subsubsection{bier --- OE bǣr}\label{bier-oe-bux1e3r-1}

Derivation: \emph{*bḗrō} \textgreater{} \emph{bǣr} (regular).

\paragraph{Derivation trace}\label{derivation-trace-14}

Proto input: \emph{*bḗrō}

\begingroup
\setlength{\fboxsep}{6pt}

\noindent\fbox{%
\begin{minipage}{0.97\linewidth}
\small
\begin{tabularx}{\linewidth}{@{}>{\raggedright\arraybackslash}p{0.460\linewidth}>{\centering\arraybackslash}X>{\raggedright\arraybackslash}p{0.460\linewidth}@{}}
\begin{minipage}[t]{\linewidth}
\centering\textbf{Earlier Germanic changes}\par
\vspace{0.35em}
\raggedright
\centering\textbf{West Germanic}\par
\raggedright
\vspace{0.2em}
\raggedright [no change]\par
\vspace{0.6em}
\centering\textbf{Northwest Germanic}\par
\raggedright
\vspace{0.2em}
\begin{tabular}{@{}>{\raggedright\arraybackslash}p{0.74\linewidth}@{\hspace{0.55em}}>{\raggedright\arraybackslash}p{0.16\linewidth}@{\hspace{0.25em}}}
\mbox{NWGmc Final Long O Raising} & \emph{*bḗru} \\
\mbox{NWGmc Long E Lowering} & \emph{*bǣru} \\
\end{tabular}
\end{minipage}
&
&
\begin{minipage}[t]{\linewidth}
\centering\textbf{Old English changes}\par
\vspace{0.35em}
\raggedright
\begin{tabular}{@{}>{\raggedright\arraybackslash}p{0.74\linewidth}@{\hspace{0.55em}}>{\raggedright\arraybackslash}p{0.16\linewidth}@{\hspace{0.25em}}}
\mbox{OE High Vowel Apocope} & \emph{*bǣr} \\
\end{tabular}
\end{minipage}
\\
\end{tabularx}
\end{minipage}%
} \endgroup

\paragraph{Commentary}\label{commentary-2}

Kroonen reconstructs the noun as \emph{*bērō-} f.~`bier' and cites Old
English bar, \emph{bær} among the reflexes
(\citeproc{ref-Kroonen2013}{Kroonen 2013, 717}). The selected input
\emph{*bḗrō} is the same lexeme in the accent notation used here.

Clark Hall and Bosworth-Toller lemmatize the noun as \emph{bær}, and
Kroonen also records \emph{bar} beside it
(\citeproc{ref-ClarkHall1960}{Clark Hall 1960};
\citeproc{ref-BosworthToller1898}{Bosworth and Toller 1898, 73};
\citeproc{ref-Kroonen2013}{Kroonen 2013, 717}). The target \emph{bǣr} is
therefore a normalized long-vowel spelling of the same noun.

From \emph{*bḗrō}, Northwest Germanic final long \emph{ō} raises to
\emph{u}, long \emph{ē} lowers to \emph{ǣ}, and high-vowel apocope
yields \emph{bǣr}. The resulting noun matches the normalized Old English
target.

\paragraph{Source note}\label{source-note-1}

Lexicographic spellings vary between \emph{bær} and \emph{bar}. The
normalized target \emph{bǣr} simply marks the same long vowel explicitly
(\citeproc{ref-ClarkHall1960}{Clark Hall 1960};
\citeproc{ref-BosworthToller1898}{Bosworth and Toller 1898, 73};
\citeproc{ref-Kroonen2013}{Kroonen 2013, 717}).

\subsubsection{both --- OE bū}\label{both-oe-bux16b-1}

Derivation: \emph{*bō} \textgreater{} \emph{bū} (regular).

\paragraph{Derivation trace}\label{derivation-trace-15}

Proto input: \emph{*bō}

\begingroup
\setlength{\fboxsep}{6pt}

\noindent\fbox{%
\begin{minipage}{0.97\linewidth}
\small
\begin{tabularx}{\linewidth}{@{}>{\raggedright\arraybackslash}p{0.540\linewidth}>{\centering\arraybackslash}X>{\raggedright\arraybackslash}p{0.300\linewidth}@{}}
\begin{minipage}[t]{\linewidth}
\centering\textbf{Earlier Germanic changes}\par
\vspace{0.35em}
\raggedright
\centering\textbf{West Germanic}\par
\raggedright
\vspace{0.2em}
\raggedright [no change]\par
\vspace{0.6em}
\centering\textbf{Northwest Germanic}\par
\raggedright
\vspace{0.2em}
\begin{tabular}{@{}>{\raggedright\arraybackslash}p{0.74\linewidth}@{\hspace{0.55em}}>{\raggedright\arraybackslash}p{0.16\linewidth}@{\hspace{0.25em}}}
NWGmc Stressed Monosyllable O Raising & \emph{*bū} \\
\end{tabular}
\end{minipage}
&
&
\begin{minipage}[t]{\linewidth}
\centering\textbf{Old English changes}\par
\vspace{0.35em}
\raggedright
\raggedright [no change]\par
\end{minipage}
\\
\end{tabularx}
\end{minipage}%
} \endgroup

\paragraph{Commentary}\label{commentary-3}

Kroonen treats the Germanic numeral under \emph{*ba-} and gives the
inherited paradigm \emph{*bai}, \emph{*bans}, \emph{*bōz}/\emph{*bōns},
\emph{*bō}, with Old English \emph{bēġen}, \emph{bā}, and neuter
\emph{bū} (\citeproc{ref-Kroonen2013}{Kroonen 2013, 47}). For the
present entry, the relevant inherited form is the unextended neuter dual
\emph{*bō}.

The older explanation of \emph{bēġen} derives it from \emph{*bō-jen-},
and Orel still gives OE \emph{bezen} (\textless{} \emph{*bō-jenō)}
beside ON \emph{báðir}, OFris \emph{bēthe}, \emph{OS be-thia}, and OHG
\emph{bēde}, beide (\citeproc{ref-Orel2003}{Orel 2003, 65}). Fulk
reports that explanation cautiously and notes Seebold's preference for a
\emph{*bō-þ-} analysis instead (\citeproc{ref-Fulk2018}{Fulk 2018, sect.
10.1}). That debate matters for \emph{bēġen} and for the extended forms
behind Modern English \emph{both}, German \emph{beide}, and Dutch
\emph{beide}; it does not displace the inherited neuter \emph{*bō}
\textgreater{} \emph{bū} treated here.

The Old English dual paradigm is well established. Brunner gives
masculine \emph{bēġen}, feminine \emph{bā}, and neuter \emph{bū} beside
\emph{bā}, with compounds such as \emph{bā} \emph{twā}, \emph{bū}
\emph{tū}, and \emph{bām} \emph{twām}
(\citeproc{ref-SieversBrunner1965}{Sievers 1965, sect. 324} Anm. 2).
Campbell and Fulk present the same basic pattern: masculine
\emph{bēġen}, feminine \emph{bā}, neuter \emph{bā}, \emph{bū}, genitive
\emph{bēġra}, \emph{bēġ(e)a}, and dative \emph{bǣm}
(\citeproc{ref-Campbell1959}{Campbell 1959, sect. 683};
\citeproc{ref-Fulk2018}{Fulk 2018, sect. 10.1}).

\emph{bū} is therefore an attested neuter dual form, not a
reconstruction. It is the cleanest target for this entry because
\emph{bēġen} belongs to the historically more contested \emph{*bō-jen-}
/ analogical zone, while \emph{bā} remains a partner form within the
dual paradigm rather than the most straightforward monosyllabic
comparison.

\emph{*bō} is a stressed monosyllabic form. Campbell cites \emph{cū},
\emph{hū}, \emph{tū}, and \emph{bū} as examples of final accented
\emph{ō} \textgreater{} \emph{ū} in the West Germanic stage leading to
Old English (\citeproc{ref-Campbell1959}{Campbell 1959, sect. 122}).
Brunner states the same development more directly: Auslautendes \emph{ō}
erscheint als û in \emph{bū} \ldots{} cu \ldots{} \emph{hū}, \emph{tū}
(\citeproc{ref-SieversBrunner1965}{Sievers 1965, sect. 69}).

The development is therefore straightforward: \emph{*bō} \textgreater{}
\emph{bū}.

\paragraph{Comparison}\label{comparison}

The comparison below is manual. It separates the inherited OE target
from the other forms that belong to the same broader lexical history.

{\def\LTcaptype{none} % do not increment counter
\begin{longtable}[]{@{}
  >{\raggedright\arraybackslash}p{(\linewidth - 6\tabcolsep) * \real{0.2500}}
  >{\raggedright\arraybackslash}p{(\linewidth - 6\tabcolsep) * \real{0.2500}}
  >{\raggedright\arraybackslash}p{(\linewidth - 6\tabcolsep) * \real{0.2500}}
  >{\raggedright\arraybackslash}p{(\linewidth - 6\tabcolsep) * \real{0.2500}}@{}}
\toprule\noalign{}
\begin{minipage}[b]{\linewidth}\raggedright
Form
\end{minipage} & \begin{minipage}[b]{\linewidth}\raggedright
Source / stage
\end{minipage} & \begin{minipage}[b]{\linewidth}\raggedright
Status
\end{minipage} & \begin{minipage}[b]{\linewidth}\raggedright
Relevance to this entry
\end{minipage} \\
\midrule\noalign{}
\endhead
\bottomrule\noalign{}
\endlastfoot
\emph{*bō} \textgreater{} \emph{bū} & PGmc neuter dual \textgreater{} OE
neuter dual & selected regular comparison & main line of the entry \\
\emph{bēġen} & OE masculine dual & attested, but historically contested
and at least partly analogical in Kroonen & real OE evidence, not the
selected target \\
\emph{bā} & OE feminine dual; also neuter variant & attested partner
form & part of the OE paradigm, but not the chosen monosyllabic
comparator \\
\emph{báðir}, \emph{beide}, \emph{both} & Norse, continental West
Germanic, Modern English extended forms & related but different
formation & useful background, not the direct continuation of OE
\emph{bū} \\
\end{longtable}
}

\subsubsection{bow --- OE bīeġan}\label{bow-oe-bux12beux121an-1}

Derivation: \emph{*báugijaną} \textgreater{} \emph{bīeġan} (regular).

\paragraph{Derivation trace}\label{derivation-trace-16}

Proto input: \emph{*báugijaną}

\begingroup
\setlength{\fboxsep}{6pt}

\noindent\fbox{%
\begin{minipage}{0.97\linewidth}
\footnotesize
\begin{tabularx}{\linewidth}{@{}>{\raggedright\arraybackslash}p{0.280\linewidth}>{\centering\arraybackslash}X>{\raggedright\arraybackslash}p{0.560\linewidth}@{}}
\begin{minipage}[t]{\linewidth}
\centering\textbf{Earlier Germanic changes}\par
\vspace{0.35em}
\raggedright
\centering\textbf{West Germanic}\par
\raggedright
\vspace{0.2em}
\raggedright [no change]\par
\vspace{0.6em}
\centering\textbf{Northwest Germanic}\par
\raggedright
\vspace{0.2em}
\raggedright [no change]\par
\end{minipage}
&
&
\begin{minipage}[t]{\linewidth}
\centering\textbf{Old English changes}\par
\vspace{0.35em}
\raggedright
\begin{tabular}{@{}>{\raggedright\arraybackslash}p{0.62\linewidth}@{\hspace{0.55em}}>{\raggedright\arraybackslash}p{0.28\linewidth}@{\hspace{0.25em}}}
OE Au Fronting & \emph{*báeugijaną} \\
OE Diphthong Leveling & \emph{*bēagijaną} \\
OE Heavy Syllable Nasal Apocope & \emph{*bēagijan} \\
OE Secondary Nasalization & \emph{*bēagijąn} \\
Sievers Law Syncope & \emph{*bēagjąn} \\
OE Velar Palatalization & \emph{*bēaʤjąn} \\
OE I Umlaut & \emph{*bīeʤjąn} \\
OE Weak Tail Reduction & \emph{*bīeʤjan} \\
OE J Loss After Heavy & \emph{*bīeʤan} \\
\end{tabular}
\end{minipage}
\\
\end{tabularx}
\end{minipage}%
} \endgroup

\paragraph{Commentary}\label{commentary-4}

Kroonen reconstructs the weak verb as \emph{*baugjan-} `to (make) bend'
and cites Old English \emph{biegan} among its reflexes
(\citeproc{ref-Kroonen2013}{Kroonen 2013}). Ringe and Taylor give the
northwest Germanic and West Saxon development more fully as PNWGmce
\emph{*baugijana} \textgreater{} \emph{*béagjan} \textgreater{} WS OE
\emph{biegan} (\citeproc{ref-RingeTaylor2014}{Ringe and Taylor 2014}).
The entry therefore concerns the weak causative member of the
bend-family, alongside the related strong verb and noun.

Clark Hall lemmatizes \emph{biegan}, and Bosworth-Toller records
\emph{bigan} with examples such as Ic \emph{bēge} \emph{mīne}
\emph{cneówa} and Se ord \emph{bīgde} upp \emph{tō} \emph{þām} hiltum
(\citeproc{ref-ClarkHall1960}{Clark Hall 1960};
\citeproc{ref-BosworthToller1898}{Bosworth and Toller 1898, 102}). The
form \emph{bīeġan} used here is a normalized spelling of that attested
Old English weak verb.

From \emph{*báugijaną}, the stem reaches pre-Old-English
\emph{*bēagjan}, after which palatalization of \emph{*gj} and i-umlaut
yield West Saxon \emph{biegan}; Campbell lists \emph{biegan} among the
regular \emph{ie} outcomes of \emph{*éa} under i-umlaut
(\citeproc{ref-RingeTaylor2014}{Ringe and Taylor 2014};
\citeproc{ref-Campbell1959}{Campbell 1959, 80}). The development is
therefore straightforward: \emph{*báugijaną} \textgreater{}
\emph{bīeġan}.

\subsubsection{fare --- OE faran}\label{fare-oe-faran-1}

Derivation: \emph{*fáraną} \textgreater{} \emph{faran} (regular).

\paragraph{Derivation trace}\label{derivation-trace-17}

Proto input: \emph{*fáraną}

\begingroup
\setlength{\fboxsep}{6pt}

\noindent\fbox{%
\begin{minipage}{0.97\linewidth}
\small
\begin{tabularx}{\linewidth}{@{}>{\raggedright\arraybackslash}p{0.300\linewidth}>{\centering\arraybackslash}X>{\raggedright\arraybackslash}p{0.540\linewidth}@{}}
\begin{minipage}[t]{\linewidth}
\centering\textbf{Earlier Germanic changes}\par
\vspace{0.35em}
\raggedright
\centering\textbf{West Germanic}\par
\raggedright
\vspace{0.2em}
\raggedright [no change]\par
\vspace{0.6em}
\centering\textbf{Northwest Germanic}\par
\raggedright
\vspace{0.2em}
\raggedright [no change]\par
\end{minipage}
&
&
\begin{minipage}[t]{\linewidth}
\centering\textbf{Old English changes}\par
\vspace{0.35em}
\raggedright
\begin{tabular}{@{}>{\raggedright\arraybackslash}p{0.74\linewidth}@{\hspace{0.55em}}>{\raggedright\arraybackslash}p{0.16\linewidth}@{\hspace{0.25em}}}
Anglo Frisian Brightening & \emph{*færaną} \\
OE A Restoration & \emph{*faraną} \\
OE Heavy Syllable Nasal Apocope & \emph{*faran} \\
OE Secondary Nasalization & \emph{*farąn} \\
OE Weak Tail Reduction & \emph{*faran} \\
\end{tabular}
\end{minipage}
\\
\end{tabularx}
\end{minipage}%
} \endgroup

\paragraph{Commentary}\label{commentary-5}

Kroonen gives the inherited strong verb as \emph{*faran-}, and Orel
gives the same lexeme as \emph{*faranan}, both with Old English
\emph{faran} among the reflexes (\citeproc{ref-Kroonen2013}{Kroonen
2013}; \citeproc{ref-Orel2003}{Orel 2003, 132}). Campbell also uses
\emph{faran} as a standard example of Old English A-restoration
(\citeproc{ref-Campbell1959}{Campbell 1959, 61}).

Clark Hall lemmatizes the strong verb as \emph{faran} and separately
records weak \emph{færan} `to frighten'; \emph{fære}, \emph{færst}, and
\emph{færð} belong to present-tense forms of \emph{faran} rather than to
the infinitive itself (\citeproc{ref-ClarkHall1960}{Clark Hall 1960}).
Bosworth-Toller preserves the same distinction
(\citeproc{ref-BosworthToller1898}{Bosworth and Toller 1898, 108}). The
selected target is therefore the attested citation infinitive
\emph{faran}.

From \emph{*fáraną}, Anglo-Frisian brightening first gives
\emph{*færaną}, but A-restoration before single \emph{r} returns
\emph{*faraną}; later apocope and weak-tail reduction yield \emph{faran}
(\citeproc{ref-Campbell1959}{Campbell 1959, 61}). Fulk's contrast with
participial faren- \textless{} \emph{*faræn-} \textless{} \emph{*faran-}
shows why fronting elsewhere in the paradigm does not alter the
infinitive headword (\citeproc{ref-Fulk2018}{Fulk 2018}).

\subsubsection{guest --- OE ġiest}\label{guest-oe-ux121iest-1}

Derivation: \emph{*gástiz} \textgreater{} \emph{ġiest} (regular).

\paragraph{Derivation trace}\label{derivation-trace-18}

Proto input: \emph{*gástiz}

\begingroup
\setlength{\fboxsep}{6pt}

\noindent\fbox{%
\begin{minipage}{0.97\linewidth}
\small
\begin{tabularx}{\linewidth}{@{}>{\raggedright\arraybackslash}p{0.460\linewidth}>{\centering\arraybackslash}X>{\raggedright\arraybackslash}p{0.460\linewidth}@{}}
\begin{minipage}[t]{\linewidth}
\centering\textbf{Earlier Germanic changes}\par
\vspace{0.35em}
\raggedright
\centering\textbf{West Germanic}\par
\raggedright
\vspace{0.2em}
\raggedright [no change]\par
\vspace{0.6em}
\centering\textbf{Northwest Germanic}\par
\raggedright
\vspace{0.2em}
\begin{tabular}{@{}>{\raggedright\arraybackslash}p{0.74\linewidth}@{\hspace{0.55em}}>{\raggedright\arraybackslash}p{0.16\linewidth}@{\hspace{0.25em}}}
\mbox{PGmc Final Z Deletion} & \emph{*gásti} \\
\end{tabular}
\end{minipage}
&
&
\begin{minipage}[t]{\linewidth}
\centering\textbf{Old English changes}\par
\vspace{0.35em}
\raggedright
\begin{tabular}{@{}>{\raggedright\arraybackslash}p{0.74\linewidth}@{\hspace{0.55em}}>{\raggedright\arraybackslash}p{0.16\linewidth}@{\hspace{0.25em}}}
Anglo Frisian Brightening & \emph{*gæsti} \\
OE Velar Palatalization & \emph{*ʤæsti} \\
OE I Umlaut & \emph{*ʤesti} \\
OE Ws Palatal Diphthongization & \emph{*ʤiesti} \\
\mbox{OE High Vowel Apocope} & \emph{*ʤiest} \\
\end{tabular}
\end{minipage}
\\
\end{tabularx}
\end{minipage}%
} \endgroup

\paragraph{Commentary}\label{commentary-6}

Campbell and Ringe-Taylor treat the noun as an ordinary i-stem whose
West Saxon development shows palatal diphthongization, while
non-West-Saxon evidence preserves forms of the \emph{gest} type
(\citeproc{ref-Campbell1959}{Campbell 1959};
\citeproc{ref-RingeTaylor2014}{Ringe and Taylor 2014}).

Bosworth-Toller and Clark Hall record the word under forms such as
\emph{gist}, \emph{gest}, \emph{giest}, and \emph{gyst}
(\citeproc{ref-BosworthToller1898}{Bosworth and Toller 1898};
\citeproc{ref-ClarkHall1960}{Clark Hall 1960}). The selected target
\emph{ġiest} is the normalized West Saxon form within that attested
family.

From \emph{*gástiz}, Anglo-Frisian brightening gives a \emph{gæst-}
stage, and i-mutation affects the front vowel before the lost
high-vocalic ending. In West Saxon the initial palatal environment then
produces \emph{ie}, so the regular outcome is \emph{ġiest}
(\citeproc{ref-Campbell1959}{Campbell 1959};
\citeproc{ref-RingeTaylor2014}{Ringe and Taylor 2014}).

\paragraph{Dialect note}\label{dialect-note-1}

West Saxon \emph{ġiest} is the selected target here. Anglian \emph{gest}
and related spellings remain real Old English comparators rather than
corrections to that choice (\citeproc{ref-RingeTaylor2014}{Ringe and
Taylor 2014}; \citeproc{ref-BosworthToller1898}{Bosworth and Toller
1898}).

\subsubsection{learn --- OE liornian}\label{learn-oe-liornian-1}

Derivation: \emph{*líznōjaną} \textgreater{} \emph{liornian} (regular).

\paragraph{Derivation trace}\label{derivation-trace-19}

Proto input: \emph{*líznōjaną}

\begingroup
\setlength{\fboxsep}{6pt}

\noindent\fbox{%
\begin{minipage}{0.97\linewidth}
\footnotesize
\begin{tabularx}{\linewidth}{@{}>{\raggedright\arraybackslash}p{0.280\linewidth}>{\centering\arraybackslash}X>{\raggedright\arraybackslash}p{0.560\linewidth}@{}}
\begin{minipage}[t]{\linewidth}
\centering\textbf{Earlier Germanic changes}\par
\vspace{0.35em}
\raggedright
\centering\textbf{West Germanic}\par
\raggedright
\vspace{0.2em}
\raggedright [no change]\par
\vspace{0.6em}
\centering\textbf{Northwest Germanic}\par
\raggedright
\vspace{0.2em}
\raggedright [no change]\par
\end{minipage}
&
&
\begin{minipage}[t]{\linewidth}
\centering\textbf{Old English changes}\par
\vspace{0.35em}
\raggedright
\begin{tabular}{@{}>{\raggedright\arraybackslash}p{0.62\linewidth}@{\hspace{0.55em}}>{\raggedright\arraybackslash}p{0.28\linewidth}@{\hspace{0.25em}}}
OE Breaking & \emph{*líornōjaną} \\
OE Heavy Syllable Nasal Apocope & \emph{*líornōjan} \\
OE Secondary Nasalization & \emph{*líornōjąn} \\
OE I Umlaut & \emph{*líornējąn} \\
OE Unstressed Long Vowel Shortening & \emph{*líornejąn} \\
OE Weak Tail Reduction & \emph{*líornejan} \\
OE Intervocalic J Vocalization & \emph{*líorneian} \\
OE Unstressed EI Contraction & \emph{*líornian} \\
\end{tabular}
\end{minipage}
\\
\end{tabularx}
\end{minipage}%
} \endgroup

\paragraph{Commentary}\label{commentary-7}

Ringe and Taylor place the verb in a class-II weak family of the
\emph{*liznō-} type (\citeproc{ref-RingeTaylor2014}{Ringe and Taylor
2014}), and Kroonen keeps the same comparative base for the Germanic
lexeme (\citeproc{ref-Kroonen2013}{Kroonen 2013}). The selected input
therefore requires no change of stem class or paradigm cell.

The non-obvious issue is dialectal. Campbell records Northumbrian
\emph{liornian} beside \emph{leornian}, and he states that the \emph{eo}
of \emph{leornian} is secondary, while Northumbrian preserves forms with
\emph{io}, where original \emph{eo} and \emph{io} remain distinct
(\citeproc{ref-Campbell1959}{Campbell 1959, sect. 123} n.~2).

The form modeled here is \emph{liornian}, the Northumbrian member of the
Old English family. Dictionary practice more often privileges
\emph{leornian} as the ordinary headword
(\citeproc{ref-ClarkHall1960}{Clark Hall 1960};
\citeproc{ref-BrightCassidyRingler1971}{Bright 1971}), but Campbell's
dialect evidence shows that \emph{liornian} is a genuine Old English
form (\citeproc{ref-Campbell1959}{Campbell 1959, sect. 123} n.~2).

This entry therefore remains compact. The point is to state clearly that
the selected target belongs to the Northumbrian side of the OE evidence
rather than to the leveled \emph{leornian} headword tradition.

From \emph{*líznōjaną}, the expected Old English developments include
rhotacism of \emph{z}, followed by breaking before \emph{r} plus
consonant, and the ordinary reduction of the weak verbal ending. The
result is \emph{liornian}, preserving the \emph{io} spelling that
Campbell associates with Northumbrian where original \emph{eo} and
\emph{io} remain distinct (\citeproc{ref-Campbell1959}{Campbell 1959,
sect. 123} n.~2).

\emph{Leornian} reflects the later \emph{eo} development that Campbell
treats as secondary {[}Campbell (\citeproc{ref-Campbell1959}{1959}),
§123 n.~2; §296{]}. The selected Northumbrian target is therefore the
regular comparison form for the \emph{i}-grade member of the family.

\paragraph{Comparison}\label{comparison-1}

The comparison below is manual. It distinguishes the regular
Northumbrian form modeled here from the better-known West Saxon
headword.

{\def\LTcaptype{none} % do not increment counter
\begin{longtable}[]{@{}
  >{\raggedright\arraybackslash}p{(\linewidth - 4\tabcolsep) * \real{0.3333}}
  >{\raggedright\arraybackslash}p{(\linewidth - 4\tabcolsep) * \real{0.3333}}
  >{\raggedright\arraybackslash}p{(\linewidth - 4\tabcolsep) * \real{0.3333}}@{}}
\toprule\noalign{}
\begin{minipage}[b]{\linewidth}\raggedright
Form
\end{minipage} & \begin{minipage}[b]{\linewidth}\raggedright
Status
\end{minipage} & \begin{minipage}[b]{\linewidth}\raggedright
Relevance to this entry
\end{minipage} \\
\midrule\noalign{}
\endhead
\bottomrule\noalign{}
\endlastfoot
\emph{*líznōjaną} -\textgreater{} liornian & computed regular output;
attested Northumbrian comparison form & selected comparison \\
\emph{leornian} & attested later \emph{eo} form and dictionary headword
& useful control, but not the target of this entry \\
\end{longtable}
}

\subsubsection{linden --- OE lind}\label{linden-oe-lind-1}

Derivation: \emph{*líndō} \textgreater{} \emph{lind} (regular).

\paragraph{Derivation trace}\label{derivation-trace-20}

Proto input: \emph{*líndō}

\begingroup
\setlength{\fboxsep}{6pt}

\noindent\fbox{%
\begin{minipage}{0.97\linewidth}
\small
\begin{tabularx}{\linewidth}{@{}>{\raggedright\arraybackslash}p{0.460\linewidth}>{\centering\arraybackslash}X>{\raggedright\arraybackslash}p{0.460\linewidth}@{}}
\begin{minipage}[t]{\linewidth}
\centering\textbf{Earlier Germanic changes}\par
\vspace{0.35em}
\raggedright
\centering\textbf{West Germanic}\par
\raggedright
\vspace{0.2em}
\raggedright [no change]\par
\vspace{0.6em}
\centering\textbf{Northwest Germanic}\par
\raggedright
\vspace{0.2em}
\begin{tabular}{@{}>{\raggedright\arraybackslash}p{0.74\linewidth}@{\hspace{0.55em}}>{\raggedright\arraybackslash}p{0.16\linewidth}@{\hspace{0.25em}}}
\mbox{NWGmc Final Long O Raising} & \emph{*líndu} \\
\end{tabular}
\end{minipage}
&
&
\begin{minipage}[t]{\linewidth}
\centering\textbf{Old English changes}\par
\vspace{0.35em}
\raggedright
\begin{tabular}{@{}>{\raggedright\arraybackslash}p{0.74\linewidth}@{\hspace{0.55em}}>{\raggedright\arraybackslash}p{0.16\linewidth}@{\hspace{0.25em}}}
\mbox{OE High Vowel Apocope} & \emph{*línd} \\
\end{tabular}
\end{minipage}
\\
\end{tabularx}
\end{minipage}%
} \endgroup

\paragraph{Commentary}\label{commentary-8}

Kroonen cites \emph{*lindō-} `lime tree' and gives Old English
\emph{lind} as the relevant reflex (\citeproc{ref-Kroonen2013}{Kroonen
2013}).

Clark Hall and Bosworth-Toller both record \emph{lind} as the noun
`lime-tree, linden' (\citeproc{ref-ClarkHall1960}{Clark Hall 1960};
\citeproc{ref-BosworthToller1898}{Bosworth and Toller 1898, 630}).

From \emph{*líndō}, Northwest Germanic final \emph{*ō} raising first
gives a \emph{*líndu} stage, and later high-vowel apocope yields
\emph{lind}. The development is therefore straightforward and regular.

\paragraph{Form note}\label{form-note-3}

The Old English noun represented here is \emph{lind}. Clark Hall also
has a separate adjectival \emph{linden} `made of linden-wood', but that
is not the noun counterpart for this entry
(\citeproc{ref-ClarkHall1960}{Clark Hall 1960}).

\subsubsection{mother --- OE mōder}\label{mother-oe-mux14dder-1}

Derivation: \emph{*mōdēr} \textgreater{} \emph{mōder} (regular).

\paragraph{Derivation trace}\label{derivation-trace-21}

Proto input: \emph{*mōdēr}

\begingroup
\setlength{\fboxsep}{6pt}

\noindent\fbox{%
\begin{minipage}{0.97\linewidth}
\small
\begin{tabularx}{\linewidth}{@{}>{\raggedright\arraybackslash}p{0.460\linewidth}>{\centering\arraybackslash}X>{\raggedright\arraybackslash}p{0.460\linewidth}@{}}
\begin{minipage}[t]{\linewidth}
\centering\textbf{Earlier Germanic changes}\par
\vspace{0.35em}
\raggedright
\centering\textbf{West Germanic}\par
\raggedright
\vspace{0.2em}
\raggedright [no change]\par
\vspace{0.6em}
\centering\textbf{Northwest Germanic}\par
\raggedright
\vspace{0.2em}
\begin{tabular}{@{}>{\raggedright\arraybackslash}p{0.74\linewidth}@{\hspace{0.55em}}>{\raggedright\arraybackslash}p{0.16\linewidth}@{\hspace{0.25em}}}
\mbox{NWGmc Long E Lowering} & \emph{*mōdǣr} \\
\end{tabular}
\end{minipage}
&
&
\begin{minipage}[t]{\linewidth}
\centering\textbf{Old English changes}\par
\vspace{0.35em}
\raggedright
\begin{tabular}{@{}>{\raggedright\arraybackslash}p{0.74\linewidth}@{\hspace{0.55em}}>{\raggedright\arraybackslash}p{0.16\linewidth}@{\hspace{0.25em}}}
OE Unstressed Long Vowel Shortening & \emph{*mōdær} \\
OE Unstressed AE Merger & \emph{*mōder} \\
\end{tabular}
\end{minipage}
\\
\end{tabularx}
\end{minipage}%
} \endgroup

\paragraph{Commentary}\label{commentary-9}

Kroonen and Orel cite the Proto-Germanic r-stem kinship noun as
\emph{*mōder-} / \emph{*mōdēr} (\citeproc{ref-Kroonen2013}{Kroonen
2013}; \citeproc{ref-Orel2003}{Orel 2003}).

The transmitted Old English headword tradition is \emph{mōdor} / modor,
with oblique \emph{mēder} in the paradigm. Clark Hall, Campbell, and
Ringe and Taylor all preserve that contrast
(\citeproc{ref-ClarkHall1960}{Clark Hall 1960};
\citeproc{ref-Campbell1959}{Campbell 1959};
\citeproc{ref-RingeTaylor2014}{Ringe and Taylor 2014}).

From \emph{*mōdēr}, the regular suffixal development yields
\emph{mōder}. That regular nominative reflex is the form represented
here, while the more familiar citation form \emph{mōdor} reflects later
levelling within the r-stem paradigm
(\citeproc{ref-Campbell1959}{Campbell 1959};
\citeproc{ref-RingeTaylor2014}{Ringe and Taylor 2014}).

\paragraph{Comparison}\label{comparison-2}

The note therefore concerns inherited vocalism rather than a different
lexeme: \emph{mōder} is the regularized nominative represented here, but
dictionaries usually print \emph{mōdor} / modor, and the oblique
evidence survives in \emph{mēder} (\citeproc{ref-ClarkHall1960}{Clark
Hall 1960}; \citeproc{ref-RingeTaylor2014}{Ringe and Taylor 2014}).

\subsubsection{show --- OE sċēawian}\label{show-oe-sux10bux113awian-1}

Derivation: \emph{*skáwōjaną} \textgreater{} \emph{sċēawian} (regular).

\paragraph{Derivation trace}\label{derivation-trace-22}

Proto input: \emph{*skáwōjaną}

\begingroup
\setlength{\fboxsep}{6pt}

\noindent\fbox{%
\begin{minipage}{0.97\linewidth}
\footnotesize
\begin{tabularx}{\linewidth}{@{}>{\raggedright\arraybackslash}p{0.280\linewidth}>{\centering\arraybackslash}X>{\raggedright\arraybackslash}p{0.560\linewidth}@{}}
\begin{minipage}[t]{\linewidth}
\centering\textbf{Earlier Germanic changes}\par
\vspace{0.35em}
\raggedright
\centering\textbf{West Germanic}\par
\raggedright
\vspace{0.2em}
\raggedright [no change]\par
\vspace{0.6em}
\centering\textbf{Northwest Germanic}\par
\raggedright
\vspace{0.2em}
\raggedright [no change]\par
\end{minipage}
&
&
\begin{minipage}[t]{\linewidth}
\centering\textbf{Old English changes}\par
\vspace{0.35em}
\raggedright
\begin{tabular}{@{}>{\raggedright\arraybackslash}p{0.62\linewidth}@{\hspace{0.55em}}>{\raggedright\arraybackslash}p{0.28\linewidth}@{\hspace{0.25em}}}
OE Aw Long Diphthong & \emph{*skḗawōjaną} \\
OE Heavy Syllable Nasal Apocope & \emph{*skḗawōjan} \\
OE Secondary Nasalization & \emph{*skḗawōjąn} \\
OE Sk Palatalization & \emph{*ʃḗawōjąn} \\
OE I Umlaut & \emph{*ʃḗawējąn} \\
OE Unstressed Long Vowel Shortening & \emph{*ʃḗawejąn} \\
OE Weak Tail Reduction & \emph{*ʃḗawejan} \\
OE Intervocalic J Vocalization & \emph{*ʃḗaweian} \\
OE Unstressed EI Contraction & \emph{*ʃḗawian} \\
\end{tabular}
\end{minipage}
\\
\end{tabularx}
\end{minipage}%
} \endgroup

\paragraph{Commentary}\label{commentary-10}

Orel and Kroonen cite a Class II verb of the type \emph{*skawōjan-},
with OE \emph{scēawian} among the reflexes (\citeproc{ref-Orel2003}{Orel
2003}; \citeproc{ref-Kroonen2013}{Kroonen 2013, 482}). Brunner likewise
records the Old English family as \emph{scēawian}, \emph{scāwian}, which
places this entry in the ordinary show-verb set rather than in a special
finite-cell workaround (\citeproc{ref-SieversBrunner1965}{Sievers
1965}).

Bright lists \emph{scēawian} (W. II.) and also the related form
\emph{scēawa} (\citeproc{ref-BrightCassidyRingler1971}{Bright 1971}).
The source tradition therefore uses \emph{scēawian}, while the target
represented here is the normalized project spelling \emph{sċēawian}.

From \emph{*skáwōjaną}, Old English \emph{aw} before a following vowel
yields \emph{ēaw}, and the Class II suffix keeps \emph{*ō} between
\emph{*w} and \emph{*j}. The development therefore runs regularly to
\emph{sċēawian}, without the direct \emph{*aw+j} problem seen in other
verb types (\citeproc{ref-Campbell1959}{Campbell 1959};
\citeproc{ref-Orel2003}{Orel 2003}).

\paragraph{Form note}\label{form-note-4}

The difference between \emph{scēawian} and \emph{sċēawian} is
orthographic normalization of initial \emph{}, not a difference of
lexeme or paradigm cell (\citeproc{ref-Campbell1959}{Campbell 1959};
\citeproc{ref-Hogg1992}{Hogg 1992}).

\subsubsection{weapon --- OE wǣpn}\label{weapon-oe-wux1e3pn-1}

Derivation: \emph{*wḗpną} \textgreater{} \emph{wǣpn} (regular).

\paragraph{Derivation trace}\label{derivation-trace-23}

Proto input: \emph{*wḗpną}

\begingroup
\setlength{\fboxsep}{6pt}

\noindent\fbox{%
\begin{minipage}{0.97\linewidth}
\small
\begin{tabularx}{\linewidth}{@{}>{\raggedright\arraybackslash}p{0.460\linewidth}>{\centering\arraybackslash}X>{\raggedright\arraybackslash}p{0.460\linewidth}@{}}
\begin{minipage}[t]{\linewidth}
\centering\textbf{Earlier Germanic changes}\par
\vspace{0.35em}
\raggedright
\centering\textbf{West Germanic}\par
\raggedright
\vspace{0.2em}
\raggedright [no change]\par
\vspace{0.6em}
\centering\textbf{Northwest Germanic}\par
\raggedright
\vspace{0.2em}
\begin{tabular}{@{}>{\raggedright\arraybackslash}p{0.74\linewidth}@{\hspace{0.55em}}>{\raggedright\arraybackslash}p{0.16\linewidth}@{\hspace{0.25em}}}
\mbox{NWGmc Long E Lowering} & \emph{*wǣpną} \\
\end{tabular}
\end{minipage}
&
&
\begin{minipage}[t]{\linewidth}
\centering\textbf{Old English changes}\par
\vspace{0.35em}
\raggedright
\begin{tabular}{@{}>{\raggedright\arraybackslash}p{0.74\linewidth}@{\hspace{0.55em}}>{\raggedright\arraybackslash}p{0.16\linewidth}@{\hspace{0.25em}}}
OE Heavy Syllable Nasal Apocope & \emph{*wǣpn} \\
\end{tabular}
\end{minipage}
\\
\end{tabularx}
\end{minipage}%
} \endgroup

\paragraph{Commentary}\label{commentary-11}

Kroonen reconstructs a double-stem noun \emph{*wēbna-} \textasciitilde{}
\emph{*wēpna-} and cites OE \emph{wæpn} among its reflexes
(\citeproc{ref-Kroonen2013}{Kroonen 2013, 617}). The selected input
\emph{*wḗpną} represents the unbroken citation-form noun rather than the
later broken simplex.

Campbell's cluster-noun discussion preserves unbroken \emph{wépn} beside
broken \emph{wépen}-type forms (\citeproc{ref-Campbell1959}{Campbell
1959, 150}; \citeproc{ref-Campbell1959}{1959, 226--27}). Bright
contrasts broken nominative \emph{wǣpen}/wapen with unbroken oblique
\emph{wǣpnes}, while Clark Hall lemmatizes the noun under \emph{wapen}
and also preserves unbroken forms in compounds and related spellings
(\citeproc{ref-BrightCassidyRingler1971}{Bright 1971, 29};
\citeproc{ref-ClarkHall1960}{Clark Hall 1960, 355}).

Northwest Germanic lowering gives \emph{wǣpn}, and loss of the final
nasal vowel leaves the unbroken cluster word-finally. The selected
target is the attested unbroken form \emph{wǣpn}.

\paragraph{Form note}\label{form-note-5}

The ordinary late West Saxon simplex headword is \emph{wǣpen}, but
Campbell's noun-class discussion also preserves unbroken \emph{wépn}
beside broken \emph{wépen}-type forms
(\citeproc{ref-Campbell1959}{Campbell 1959, 150};
\citeproc{ref-Campbell1959}{1959, 226--27}). \emph{wǣpnes} remains the
regular unbroken oblique comparator
(\citeproc{ref-BrightCassidyRingler1971}{Bright 1971, 29};
\citeproc{ref-ClarkHall1960}{Clark Hall 1960, 355}).

\clearpage

\subsection{Variant C: minimal regular style
(experimental)}\label{variant-c-minimal-regular-style-experimental}

These entries keep the heading, derivation line, and boxed trace, then
retain only explicit short note sections. Manual comparison tables are
replaced by a mechanical editorial placeholder.

\subsubsection{bake --- OE bacan}\label{bake-oe-bacan-2}

Derivation: \emph{*bákaną} \textgreater{} \emph{bacan} (regular).

\paragraph{Derivation trace}\label{derivation-trace-24}

Proto input: \emph{*bákaną}

\begingroup
\setlength{\fboxsep}{6pt}

\noindent\fbox{%
\begin{minipage}{0.97\linewidth}
\small
\begin{tabularx}{\linewidth}{@{}>{\raggedright\arraybackslash}p{0.300\linewidth}>{\centering\arraybackslash}X>{\raggedright\arraybackslash}p{0.540\linewidth}@{}}
\begin{minipage}[t]{\linewidth}
\centering\textbf{Earlier Germanic changes}\par
\vspace{0.35em}
\raggedright
\centering\textbf{West Germanic}\par
\raggedright
\vspace{0.2em}
\raggedright [no change]\par
\vspace{0.6em}
\centering\textbf{Northwest Germanic}\par
\raggedright
\vspace{0.2em}
\raggedright [no change]\par
\end{minipage}
&
&
\begin{minipage}[t]{\linewidth}
\centering\textbf{Old English changes}\par
\vspace{0.35em}
\raggedright
\begin{tabular}{@{}>{\raggedright\arraybackslash}p{0.74\linewidth}@{\hspace{0.55em}}>{\raggedright\arraybackslash}p{0.16\linewidth}@{\hspace{0.25em}}}
Anglo Frisian Brightening & \emph{*bækaną} \\
OE A Restoration & \emph{*bakaną} \\
OE Heavy Syllable Nasal Apocope & \emph{*bakan} \\
OE Secondary Nasalization & \emph{*bakąn} \\
OE Weak Tail Reduction & \emph{*bakan} \\
\end{tabular}
\end{minipage}
\\
\end{tabularx}
\end{minipage}%
} \endgroup

\subsubsection{beech --- OE bōc}\label{beech-oe-bux14dc-2}

Derivation: \emph{*bōkō} \textgreater{} \emph{bōc} (regular).

\paragraph{Derivation trace}\label{derivation-trace-25}

Proto input: \emph{*bōkō}

\begingroup
\setlength{\fboxsep}{6pt}

\noindent\fbox{%
\begin{minipage}{0.97\linewidth}
\small
\begin{tabularx}{\linewidth}{@{}>{\raggedright\arraybackslash}p{0.460\linewidth}>{\centering\arraybackslash}X>{\raggedright\arraybackslash}p{0.460\linewidth}@{}}
\begin{minipage}[t]{\linewidth}
\centering\textbf{Earlier Germanic changes}\par
\vspace{0.35em}
\raggedright
\centering\textbf{West Germanic}\par
\raggedright
\vspace{0.2em}
\raggedright [no change]\par
\vspace{0.6em}
\centering\textbf{Northwest Germanic}\par
\raggedright
\vspace{0.2em}
\begin{tabular}{@{}>{\raggedright\arraybackslash}p{0.74\linewidth}@{\hspace{0.55em}}>{\raggedright\arraybackslash}p{0.16\linewidth}@{\hspace{0.25em}}}
\mbox{NWGmc Final Long O Raising} & \emph{*bōku} \\
\end{tabular}
\end{minipage}
&
&
\begin{minipage}[t]{\linewidth}
\centering\textbf{Old English changes}\par
\vspace{0.35em}
\raggedright
\begin{tabular}{@{}>{\raggedright\arraybackslash}p{0.74\linewidth}@{\hspace{0.55em}}>{\raggedright\arraybackslash}p{0.16\linewidth}@{\hspace{0.25em}}}
\mbox{OE High Vowel Apocope} & \emph{*bōk} \\
\end{tabular}
\end{minipage}
\\
\end{tabularx}
\end{minipage}%
} \endgroup

\subsubsection{bier --- OE bǣr}\label{bier-oe-bux1e3r-2}

Derivation: \emph{*bḗrō} \textgreater{} \emph{bǣr} (regular).

\paragraph{Derivation trace}\label{derivation-trace-26}

Proto input: \emph{*bḗrō}

\begingroup
\setlength{\fboxsep}{6pt}

\noindent\fbox{%
\begin{minipage}{0.97\linewidth}
\small
\begin{tabularx}{\linewidth}{@{}>{\raggedright\arraybackslash}p{0.460\linewidth}>{\centering\arraybackslash}X>{\raggedright\arraybackslash}p{0.460\linewidth}@{}}
\begin{minipage}[t]{\linewidth}
\centering\textbf{Earlier Germanic changes}\par
\vspace{0.35em}
\raggedright
\centering\textbf{West Germanic}\par
\raggedright
\vspace{0.2em}
\raggedright [no change]\par
\vspace{0.6em}
\centering\textbf{Northwest Germanic}\par
\raggedright
\vspace{0.2em}
\begin{tabular}{@{}>{\raggedright\arraybackslash}p{0.74\linewidth}@{\hspace{0.55em}}>{\raggedright\arraybackslash}p{0.16\linewidth}@{\hspace{0.25em}}}
\mbox{NWGmc Final Long O Raising} & \emph{*bḗru} \\
\mbox{NWGmc Long E Lowering} & \emph{*bǣru} \\
\end{tabular}
\end{minipage}
&
&
\begin{minipage}[t]{\linewidth}
\centering\textbf{Old English changes}\par
\vspace{0.35em}
\raggedright
\begin{tabular}{@{}>{\raggedright\arraybackslash}p{0.74\linewidth}@{\hspace{0.55em}}>{\raggedright\arraybackslash}p{0.16\linewidth}@{\hspace{0.25em}}}
\mbox{OE High Vowel Apocope} & \emph{*bǣr} \\
\end{tabular}
\end{minipage}
\\
\end{tabularx}
\end{minipage}%
} \endgroup

\paragraph{Source note}\label{source-note-2}

Lexicographic spellings vary between \emph{bær} and \emph{bar}. The
normalized target \emph{bǣr} simply marks the same long vowel explicitly
(\citeproc{ref-ClarkHall1960}{Clark Hall 1960};
\citeproc{ref-BosworthToller1898}{Bosworth and Toller 1898, 73};
\citeproc{ref-Kroonen2013}{Kroonen 2013, 717}).

\subsubsection{both --- OE bū}\label{both-oe-bux16b-2}

Derivation: \emph{*bō} \textgreater{} \emph{bū} (regular).

\paragraph{Derivation trace}\label{derivation-trace-27}

Proto input: \emph{*bō}

\begingroup
\setlength{\fboxsep}{6pt}

\noindent\fbox{%
\begin{minipage}{0.97\linewidth}
\small
\begin{tabularx}{\linewidth}{@{}>{\raggedright\arraybackslash}p{0.540\linewidth}>{\centering\arraybackslash}X>{\raggedright\arraybackslash}p{0.300\linewidth}@{}}
\begin{minipage}[t]{\linewidth}
\centering\textbf{Earlier Germanic changes}\par
\vspace{0.35em}
\raggedright
\centering\textbf{West Germanic}\par
\raggedright
\vspace{0.2em}
\raggedright [no change]\par
\vspace{0.6em}
\centering\textbf{Northwest Germanic}\par
\raggedright
\vspace{0.2em}
\begin{tabular}{@{}>{\raggedright\arraybackslash}p{0.74\linewidth}@{\hspace{0.55em}}>{\raggedright\arraybackslash}p{0.16\linewidth}@{\hspace{0.25em}}}
NWGmc Stressed Monosyllable O Raising & \emph{*bū} \\
\end{tabular}
\end{minipage}
&
&
\begin{minipage}[t]{\linewidth}
\centering\textbf{Old English changes}\par
\vspace{0.35em}
\raggedright
\raggedright [no change]\par
\end{minipage}
\\
\end{tabularx}
\end{minipage}%
} \endgroup

\paragraph{Comparison}\label{comparison-3}

\emph{Manual comparison retained only in fuller variants.}

\subsubsection{bow --- OE bīeġan}\label{bow-oe-bux12beux121an-2}

Derivation: \emph{*báugijaną} \textgreater{} \emph{bīeġan} (regular).

\paragraph{Derivation trace}\label{derivation-trace-28}

Proto input: \emph{*báugijaną}

\begingroup
\setlength{\fboxsep}{6pt}

\noindent\fbox{%
\begin{minipage}{0.97\linewidth}
\footnotesize
\begin{tabularx}{\linewidth}{@{}>{\raggedright\arraybackslash}p{0.280\linewidth}>{\centering\arraybackslash}X>{\raggedright\arraybackslash}p{0.560\linewidth}@{}}
\begin{minipage}[t]{\linewidth}
\centering\textbf{Earlier Germanic changes}\par
\vspace{0.35em}
\raggedright
\centering\textbf{West Germanic}\par
\raggedright
\vspace{0.2em}
\raggedright [no change]\par
\vspace{0.6em}
\centering\textbf{Northwest Germanic}\par
\raggedright
\vspace{0.2em}
\raggedright [no change]\par
\end{minipage}
&
&
\begin{minipage}[t]{\linewidth}
\centering\textbf{Old English changes}\par
\vspace{0.35em}
\raggedright
\begin{tabular}{@{}>{\raggedright\arraybackslash}p{0.62\linewidth}@{\hspace{0.55em}}>{\raggedright\arraybackslash}p{0.28\linewidth}@{\hspace{0.25em}}}
OE Au Fronting & \emph{*báeugijaną} \\
OE Diphthong Leveling & \emph{*bēagijaną} \\
OE Heavy Syllable Nasal Apocope & \emph{*bēagijan} \\
OE Secondary Nasalization & \emph{*bēagijąn} \\
Sievers Law Syncope & \emph{*bēagjąn} \\
OE Velar Palatalization & \emph{*bēaʤjąn} \\
OE I Umlaut & \emph{*bīeʤjąn} \\
OE Weak Tail Reduction & \emph{*bīeʤjan} \\
OE J Loss After Heavy & \emph{*bīeʤan} \\
\end{tabular}
\end{minipage}
\\
\end{tabularx}
\end{minipage}%
} \endgroup

\subsubsection{fare --- OE faran}\label{fare-oe-faran-2}

Derivation: \emph{*fáraną} \textgreater{} \emph{faran} (regular).

\paragraph{Derivation trace}\label{derivation-trace-29}

Proto input: \emph{*fáraną}

\begingroup
\setlength{\fboxsep}{6pt}

\noindent\fbox{%
\begin{minipage}{0.97\linewidth}
\small
\begin{tabularx}{\linewidth}{@{}>{\raggedright\arraybackslash}p{0.300\linewidth}>{\centering\arraybackslash}X>{\raggedright\arraybackslash}p{0.540\linewidth}@{}}
\begin{minipage}[t]{\linewidth}
\centering\textbf{Earlier Germanic changes}\par
\vspace{0.35em}
\raggedright
\centering\textbf{West Germanic}\par
\raggedright
\vspace{0.2em}
\raggedright [no change]\par
\vspace{0.6em}
\centering\textbf{Northwest Germanic}\par
\raggedright
\vspace{0.2em}
\raggedright [no change]\par
\end{minipage}
&
&
\begin{minipage}[t]{\linewidth}
\centering\textbf{Old English changes}\par
\vspace{0.35em}
\raggedright
\begin{tabular}{@{}>{\raggedright\arraybackslash}p{0.74\linewidth}@{\hspace{0.55em}}>{\raggedright\arraybackslash}p{0.16\linewidth}@{\hspace{0.25em}}}
Anglo Frisian Brightening & \emph{*færaną} \\
OE A Restoration & \emph{*faraną} \\
OE Heavy Syllable Nasal Apocope & \emph{*faran} \\
OE Secondary Nasalization & \emph{*farąn} \\
OE Weak Tail Reduction & \emph{*faran} \\
\end{tabular}
\end{minipage}
\\
\end{tabularx}
\end{minipage}%
} \endgroup

\subsubsection{guest --- OE ġiest}\label{guest-oe-ux121iest-2}

Derivation: \emph{*gástiz} \textgreater{} \emph{ġiest} (regular).

\paragraph{Derivation trace}\label{derivation-trace-30}

Proto input: \emph{*gástiz}

\begingroup
\setlength{\fboxsep}{6pt}

\noindent\fbox{%
\begin{minipage}{0.97\linewidth}
\small
\begin{tabularx}{\linewidth}{@{}>{\raggedright\arraybackslash}p{0.460\linewidth}>{\centering\arraybackslash}X>{\raggedright\arraybackslash}p{0.460\linewidth}@{}}
\begin{minipage}[t]{\linewidth}
\centering\textbf{Earlier Germanic changes}\par
\vspace{0.35em}
\raggedright
\centering\textbf{West Germanic}\par
\raggedright
\vspace{0.2em}
\raggedright [no change]\par
\vspace{0.6em}
\centering\textbf{Northwest Germanic}\par
\raggedright
\vspace{0.2em}
\begin{tabular}{@{}>{\raggedright\arraybackslash}p{0.74\linewidth}@{\hspace{0.55em}}>{\raggedright\arraybackslash}p{0.16\linewidth}@{\hspace{0.25em}}}
\mbox{PGmc Final Z Deletion} & \emph{*gásti} \\
\end{tabular}
\end{minipage}
&
&
\begin{minipage}[t]{\linewidth}
\centering\textbf{Old English changes}\par
\vspace{0.35em}
\raggedright
\begin{tabular}{@{}>{\raggedright\arraybackslash}p{0.74\linewidth}@{\hspace{0.55em}}>{\raggedright\arraybackslash}p{0.16\linewidth}@{\hspace{0.25em}}}
Anglo Frisian Brightening & \emph{*gæsti} \\
OE Velar Palatalization & \emph{*ʤæsti} \\
OE I Umlaut & \emph{*ʤesti} \\
OE Ws Palatal Diphthongization & \emph{*ʤiesti} \\
\mbox{OE High Vowel Apocope} & \emph{*ʤiest} \\
\end{tabular}
\end{minipage}
\\
\end{tabularx}
\end{minipage}%
} \endgroup

\paragraph{Dialect note}\label{dialect-note-2}

West Saxon \emph{ġiest} is the selected target here. Anglian \emph{gest}
and related spellings remain real Old English comparators rather than
corrections to that choice (\citeproc{ref-RingeTaylor2014}{Ringe and
Taylor 2014}; \citeproc{ref-BosworthToller1898}{Bosworth and Toller
1898}).

\subsubsection{learn --- OE liornian}\label{learn-oe-liornian-2}

Derivation: \emph{*líznōjaną} \textgreater{} \emph{liornian} (regular).

\paragraph{Derivation trace}\label{derivation-trace-31}

Proto input: \emph{*líznōjaną}

\begingroup
\setlength{\fboxsep}{6pt}

\noindent\fbox{%
\begin{minipage}{0.97\linewidth}
\footnotesize
\begin{tabularx}{\linewidth}{@{}>{\raggedright\arraybackslash}p{0.280\linewidth}>{\centering\arraybackslash}X>{\raggedright\arraybackslash}p{0.560\linewidth}@{}}
\begin{minipage}[t]{\linewidth}
\centering\textbf{Earlier Germanic changes}\par
\vspace{0.35em}
\raggedright
\centering\textbf{West Germanic}\par
\raggedright
\vspace{0.2em}
\raggedright [no change]\par
\vspace{0.6em}
\centering\textbf{Northwest Germanic}\par
\raggedright
\vspace{0.2em}
\raggedright [no change]\par
\end{minipage}
&
&
\begin{minipage}[t]{\linewidth}
\centering\textbf{Old English changes}\par
\vspace{0.35em}
\raggedright
\begin{tabular}{@{}>{\raggedright\arraybackslash}p{0.62\linewidth}@{\hspace{0.55em}}>{\raggedright\arraybackslash}p{0.28\linewidth}@{\hspace{0.25em}}}
OE Breaking & \emph{*líornōjaną} \\
OE Heavy Syllable Nasal Apocope & \emph{*líornōjan} \\
OE Secondary Nasalization & \emph{*líornōjąn} \\
OE I Umlaut & \emph{*líornējąn} \\
OE Unstressed Long Vowel Shortening & \emph{*líornejąn} \\
OE Weak Tail Reduction & \emph{*líornejan} \\
OE Intervocalic J Vocalization & \emph{*líorneian} \\
OE Unstressed EI Contraction & \emph{*líornian} \\
\end{tabular}
\end{minipage}
\\
\end{tabularx}
\end{minipage}%
} \endgroup

\paragraph{Comparison}\label{comparison-4}

\emph{Manual comparison retained only in fuller variants.}

\subsubsection{linden --- OE lind}\label{linden-oe-lind-2}

Derivation: \emph{*líndō} \textgreater{} \emph{lind} (regular).

\paragraph{Derivation trace}\label{derivation-trace-32}

Proto input: \emph{*líndō}

\begingroup
\setlength{\fboxsep}{6pt}

\noindent\fbox{%
\begin{minipage}{0.97\linewidth}
\small
\begin{tabularx}{\linewidth}{@{}>{\raggedright\arraybackslash}p{0.460\linewidth}>{\centering\arraybackslash}X>{\raggedright\arraybackslash}p{0.460\linewidth}@{}}
\begin{minipage}[t]{\linewidth}
\centering\textbf{Earlier Germanic changes}\par
\vspace{0.35em}
\raggedright
\centering\textbf{West Germanic}\par
\raggedright
\vspace{0.2em}
\raggedright [no change]\par
\vspace{0.6em}
\centering\textbf{Northwest Germanic}\par
\raggedright
\vspace{0.2em}
\begin{tabular}{@{}>{\raggedright\arraybackslash}p{0.74\linewidth}@{\hspace{0.55em}}>{\raggedright\arraybackslash}p{0.16\linewidth}@{\hspace{0.25em}}}
\mbox{NWGmc Final Long O Raising} & \emph{*líndu} \\
\end{tabular}
\end{minipage}
&
&
\begin{minipage}[t]{\linewidth}
\centering\textbf{Old English changes}\par
\vspace{0.35em}
\raggedright
\begin{tabular}{@{}>{\raggedright\arraybackslash}p{0.74\linewidth}@{\hspace{0.55em}}>{\raggedright\arraybackslash}p{0.16\linewidth}@{\hspace{0.25em}}}
\mbox{OE High Vowel Apocope} & \emph{*línd} \\
\end{tabular}
\end{minipage}
\\
\end{tabularx}
\end{minipage}%
} \endgroup

\paragraph{Form note}\label{form-note-6}

The Old English noun represented here is \emph{lind}. Clark Hall also
has a separate adjectival \emph{linden} `made of linden-wood', but that
is not the noun counterpart for this entry
(\citeproc{ref-ClarkHall1960}{Clark Hall 1960}).

\subsubsection{mother --- OE mōder}\label{mother-oe-mux14dder-2}

Derivation: \emph{*mōdēr} \textgreater{} \emph{mōder} (regular).

\paragraph{Derivation trace}\label{derivation-trace-33}

Proto input: \emph{*mōdēr}

\begingroup
\setlength{\fboxsep}{6pt}

\noindent\fbox{%
\begin{minipage}{0.97\linewidth}
\small
\begin{tabularx}{\linewidth}{@{}>{\raggedright\arraybackslash}p{0.460\linewidth}>{\centering\arraybackslash}X>{\raggedright\arraybackslash}p{0.460\linewidth}@{}}
\begin{minipage}[t]{\linewidth}
\centering\textbf{Earlier Germanic changes}\par
\vspace{0.35em}
\raggedright
\centering\textbf{West Germanic}\par
\raggedright
\vspace{0.2em}
\raggedright [no change]\par
\vspace{0.6em}
\centering\textbf{Northwest Germanic}\par
\raggedright
\vspace{0.2em}
\begin{tabular}{@{}>{\raggedright\arraybackslash}p{0.74\linewidth}@{\hspace{0.55em}}>{\raggedright\arraybackslash}p{0.16\linewidth}@{\hspace{0.25em}}}
\mbox{NWGmc Long E Lowering} & \emph{*mōdǣr} \\
\end{tabular}
\end{minipage}
&
&
\begin{minipage}[t]{\linewidth}
\centering\textbf{Old English changes}\par
\vspace{0.35em}
\raggedright
\begin{tabular}{@{}>{\raggedright\arraybackslash}p{0.74\linewidth}@{\hspace{0.55em}}>{\raggedright\arraybackslash}p{0.16\linewidth}@{\hspace{0.25em}}}
OE Unstressed Long Vowel Shortening & \emph{*mōdær} \\
OE Unstressed AE Merger & \emph{*mōder} \\
\end{tabular}
\end{minipage}
\\
\end{tabularx}
\end{minipage}%
} \endgroup

\paragraph{Comparison}\label{comparison-5}

The note therefore concerns inherited vocalism rather than a different
lexeme: \emph{mōder} is the regularized nominative represented here, but
dictionaries usually print \emph{mōdor} / modor, and the oblique
evidence survives in \emph{mēder} (\citeproc{ref-ClarkHall1960}{Clark
Hall 1960}; \citeproc{ref-RingeTaylor2014}{Ringe and Taylor 2014}).

\subsubsection{show --- OE sċēawian}\label{show-oe-sux10bux113awian-2}

Derivation: \emph{*skáwōjaną} \textgreater{} \emph{sċēawian} (regular).

\paragraph{Derivation trace}\label{derivation-trace-34}

Proto input: \emph{*skáwōjaną}

\begingroup
\setlength{\fboxsep}{6pt}

\noindent\fbox{%
\begin{minipage}{0.97\linewidth}
\footnotesize
\begin{tabularx}{\linewidth}{@{}>{\raggedright\arraybackslash}p{0.280\linewidth}>{\centering\arraybackslash}X>{\raggedright\arraybackslash}p{0.560\linewidth}@{}}
\begin{minipage}[t]{\linewidth}
\centering\textbf{Earlier Germanic changes}\par
\vspace{0.35em}
\raggedright
\centering\textbf{West Germanic}\par
\raggedright
\vspace{0.2em}
\raggedright [no change]\par
\vspace{0.6em}
\centering\textbf{Northwest Germanic}\par
\raggedright
\vspace{0.2em}
\raggedright [no change]\par
\end{minipage}
&
&
\begin{minipage}[t]{\linewidth}
\centering\textbf{Old English changes}\par
\vspace{0.35em}
\raggedright
\begin{tabular}{@{}>{\raggedright\arraybackslash}p{0.62\linewidth}@{\hspace{0.55em}}>{\raggedright\arraybackslash}p{0.28\linewidth}@{\hspace{0.25em}}}
OE Aw Long Diphthong & \emph{*skḗawōjaną} \\
OE Heavy Syllable Nasal Apocope & \emph{*skḗawōjan} \\
OE Secondary Nasalization & \emph{*skḗawōjąn} \\
OE Sk Palatalization & \emph{*ʃḗawōjąn} \\
OE I Umlaut & \emph{*ʃḗawējąn} \\
OE Unstressed Long Vowel Shortening & \emph{*ʃḗawejąn} \\
OE Weak Tail Reduction & \emph{*ʃḗawejan} \\
OE Intervocalic J Vocalization & \emph{*ʃḗaweian} \\
OE Unstressed EI Contraction & \emph{*ʃḗawian} \\
\end{tabular}
\end{minipage}
\\
\end{tabularx}
\end{minipage}%
} \endgroup

\paragraph{Form note}\label{form-note-7}

The difference between \emph{scēawian} and \emph{sċēawian} is
orthographic normalization of initial \emph{}, not a difference of
lexeme or paradigm cell (\citeproc{ref-Campbell1959}{Campbell 1959};
\citeproc{ref-Hogg1992}{Hogg 1992}).

\subsubsection{weapon --- OE wǣpn}\label{weapon-oe-wux1e3pn-2}

Derivation: \emph{*wḗpną} \textgreater{} \emph{wǣpn} (regular).

\paragraph{Derivation trace}\label{derivation-trace-35}

Proto input: \emph{*wḗpną}

\begingroup
\setlength{\fboxsep}{6pt}

\noindent\fbox{%
\begin{minipage}{0.97\linewidth}
\small
\begin{tabularx}{\linewidth}{@{}>{\raggedright\arraybackslash}p{0.460\linewidth}>{\centering\arraybackslash}X>{\raggedright\arraybackslash}p{0.460\linewidth}@{}}
\begin{minipage}[t]{\linewidth}
\centering\textbf{Earlier Germanic changes}\par
\vspace{0.35em}
\raggedright
\centering\textbf{West Germanic}\par
\raggedright
\vspace{0.2em}
\raggedright [no change]\par
\vspace{0.6em}
\centering\textbf{Northwest Germanic}\par
\raggedright
\vspace{0.2em}
\begin{tabular}{@{}>{\raggedright\arraybackslash}p{0.74\linewidth}@{\hspace{0.55em}}>{\raggedright\arraybackslash}p{0.16\linewidth}@{\hspace{0.25em}}}
\mbox{NWGmc Long E Lowering} & \emph{*wǣpną} \\
\end{tabular}
\end{minipage}
&
&
\begin{minipage}[t]{\linewidth}
\centering\textbf{Old English changes}\par
\vspace{0.35em}
\raggedright
\begin{tabular}{@{}>{\raggedright\arraybackslash}p{0.74\linewidth}@{\hspace{0.55em}}>{\raggedright\arraybackslash}p{0.16\linewidth}@{\hspace{0.25em}}}
OE Heavy Syllable Nasal Apocope & \emph{*wǣpn} \\
\end{tabular}
\end{minipage}
\\
\end{tabularx}
\end{minipage}%
} \endgroup

\paragraph{Form note}\label{form-note-8}

The ordinary late West Saxon simplex headword is \emph{wǣpen}, but
Campbell's noun-class discussion also preserves unbroken \emph{wépn}
beside broken \emph{wépen}-type forms
(\citeproc{ref-Campbell1959}{Campbell 1959, 150};
\citeproc{ref-Campbell1959}{1959, 226--27}). \emph{wǣpnes} remains the
regular unbroken oblique comparator
(\citeproc{ref-BrightCassidyRingler1971}{Bright 1971, 29};
\citeproc{ref-ClarkHall1960}{Clark Hall 1960, 355}).

\clearpage

\subsection*{References}\label{references}
\addcontentsline{toc}{subsection}{References}

\protect\phantomsection\label{refs}
\begin{CSLReferences}{1}{1}
\bibitem[\citeproctext]{ref-BosworthToller1898}
Bosworth, Joseph, and T. Northcote Toller. 1898. \emph{An {Anglo-Saxon}
Dictionary}. Clarendon Press.

\bibitem[\citeproctext]{ref-BrightCassidyRingler1971}
Bright, James W. 1971. \emph{Bright's {Old English} Grammar and Reader}.
3rd edn. Edited by Frederic G. Cassidy and Richard N. Ringler. Holt,
Rinehart; Winston.

\bibitem[\citeproctext]{ref-Campbell1959}
Campbell, Alistair. 1959. \emph{Old {English} Grammar}. Clarendon Press.

\bibitem[\citeproctext]{ref-ClarkHall1960}
Clark Hall, J. R. 1960. \emph{A Concise {Anglo-Saxon} Dictionary}. 4th
edn. University of Toronto Press.

\bibitem[\citeproctext]{ref-Fulk2018}
Fulk, R. D. 2018. \emph{A Comparative Grammar of the Early {Germanic}
Languages}. Vol. 3. Studies in Germanic Linguistics. John Benjamins.

\bibitem[\citeproctext]{ref-Hogg1992}
Hogg, Richard M. 1992. \emph{A Grammar of Old English. Volume 1:
Phonology}. Blackwell.

\bibitem[\citeproctext]{ref-Kroonen2013}
Kroonen, Guus. 2013. \emph{Etymological Dictionary of {Proto-Germanic}}.
Vol. 11. Leiden Indo-European Etymological Dictionary Series. Brill.

\bibitem[\citeproctext]{ref-Orel2003}
Orel, Vladimir. 2003. \emph{A Handbook of {Germanic} Etymology}. Brill.

\bibitem[\citeproctext]{ref-RingeTaylor2014}
Ringe, Don, and Ann Taylor. 2014. \emph{The Development of {Old
English}}. Vol. 2. A Linguistic History of English. Oxford University
Press.

\bibitem[\citeproctext]{ref-SieversBrunner1965}
Sievers, Eduard. 1965. \emph{Altenglische Grammatik}. 3rd edn. Edited by
Karl Brunner. Niemeyer.

\end{CSLReferences}

\end{document}
